% Paper 3: The Qameah as Clifford Gate
% Based on outline_v1.md (outputs/research/paper3/)
% Target: arXiv math-ph, cross-list quant-ph + math.CO; journal: JMP
% Focus: GL(2,Z_N) = Clifford, Proca eigenvalue, BRST code, N=11 selection, Newton sums
\pdfoutput=1
\documentclass[aip,jmp,reprint,superscriptaddress,floatfix,nofootinbib]{revtex4-2}

\usepackage[T1]{fontenc}
\usepackage{amsmath,amssymb,amsfonts}
\usepackage{amsthm}
\usepackage{graphicx}
\usepackage{hyperref}
\usepackage{xcolor}
\usepackage{booktabs}
\usepackage{enumitem}
\usepackage{braket}

% Theorem environments
\newtheorem{theorem}{Theorem}
\newtheorem{lemma}[theorem]{Lemma}
\newtheorem{corollary}[theorem]{Corollary}
\newtheorem{proposition}[theorem]{Proposition}
\theoremstyle{definition}
\newtheorem{definition}[theorem]{Definition}
\theoremstyle{remark}
\newtheorem{remark}[theorem]{Remark}
\newtheorem{conjecture}[theorem]{Conjecture}

% Custom commands
\newcommand{\ZN}{\mathbb{Z}_N}
\newcommand{\GL}{\mathrm{GL}}
\newcommand{\SL}{\mathrm{SL}}
\newcommand{\HW}{\mathrm{HW}}
\newcommand{\ie}{\textit{i.e.}}
\newcommand{\eg}{\textit{e.g.}}

\begin{document}

\title{The Qameah as Clifford Gate: Magic Squares, Quantum Error Correction,
       and Massive Gauge Fields}

\author{Hr\'{o}ar {\TH}\'{o}r Reynisson}
\thanks{ORCID: 0009-0003-6518-121X}
\affiliation{I{\dh}nt{\ae}knistofnun (IceTec), Reykjav{\'\i}k, Iceland}

\author{Claude Opus 4.6}
\thanks{AI research collaborator, Anthropic}
\affiliation{Anthropic, San Francisco, CA, USA}

\date{\today}

\begin{abstract}
We show that the $\GL(2,\ZN)$ construction underlying $N \times N$
de la Loub\`{e}re magic squares is isomorphic to the Clifford group for
$N$-dimensional quantum systems (qu$N$its) when $N$ is an odd prime.
A single algebraic parameter $s = 2^{-1} \bmod N = (N{+}1)/2$ governs
the entire structure: the Heisenberg--Weyl labeling of entries, the
circulant/Hankel decomposition $Q = S + A$, the anticommutation
$\{S,A\} = 0$, the Proca equation $B_2^T B_2\, A = N\, A$ on the
complete graph~$K_N$, and the nilpotent BRST operators $L_\pm$ with
$L_\pm^2 = 0$.  The first BRST cohomology $H^1(L_+)$ is
one-dimensional for every prime $N$ tested ($3 \le N \le 61$),
yielding a quantum code $[[N, 1, (N{+}1)/2]]$ whose distance
$d = (N{+}1)/2$ equals $s$ itself.  We prove universal Newton sum
identities for the Lehmer roots $r_k = 2\cos(2\pi k/N)$: odd power
sums equal $-4^n$ (range $p < N$) and even power sums equal
$(N\binom{2n}{n} - 4^n)/2$ (range $p < 2N$).  For $N = 11$, four
independent selection criteria converge uniquely: the
Fibonacci--triangular coincidence $\binom{11}{2} = 55 = T_{10} =
F_{10}$, the $E_8$ root count $2(N^2 - 1) = 240$, the perfect code
distance $d = 6$, and the Euler characteristic $\chi(K_{11}) = N^2 =
|\HW(N)/Z(\HW)|$.  These results connect magic square theory, discrete
gauge theory, and quantum error correction through a single algebraic
source with zero free parameters.
\end{abstract}

\maketitle

%=====================================================================
\section{Introduction}
\label{sec:intro}

For odd prime~$N$, the de la Loub\`{e}re construction of an $N \times N$
magic square is equivalent to a linear map $g \in \GL(2,\ZN)$ acting on
pairs of mutually orthogonal Latin squares (MOLS).  Setting the
parameter $s = 2^{-1} \bmod N = (N{+}1)/2$ yields a
distinguished magic square---the \emph{Qameah} of order~$N$ (the
$\GL(2,\ZN)$ magic square with parameter $s = (N{+}1)/2$; a
distinguished member of the de la Loub\`{e}re family satisfying the
circulant criterion of Proposition~\ref{prop:circulant})---whose
algebraic structure we show to be identical to a single-qu$N$it
Clifford gate.
The identification is a \emph{structural isomorphism}: the
$\GL(2,\ZN)$ gate acts on Heisenberg--Weyl operator labels by matrix
multiplication, and this action equals the conjugation action of the
corresponding Clifford unitary~\cite{Appleby2005, Gross2006}.

From this single datum~$s$, with zero free parameters, we derive five
results:
\begin{enumerate}[label=(\roman*)]
\item \textbf{Clifford gate} (Theorem~\ref{thm:clifford-gate}):
  $g$ defines a Clifford unitary $U_g$ with $\det(g) = s$
  (extended Clifford group).
\item \textbf{Proca eigenvalue} (Theorem~\ref{thm:proca}):
  $B_2^T B_2\, A = N\, A$ on~$K_N$, a massive spin-1 field with
  $m^2 = N$.
\item \textbf{BRST code} (Theorem~\ref{thm:code}):
  $L_\pm^2 = 0$ yields a quantum code $[[N, 1, (N{+}1)/2]]$
  with $d = s$.
\item \textbf{$N = 11$ selection} (Theorem~\ref{thm:euler}):
  four independent criteria converge uniquely at $N = 11$.
\item \textbf{Newton sums} (Theorem~\ref{thm:newton}):
  universal identities for Lehmer roots $r_k = 2\cos(2\pi k/N)$.
\end{enumerate}
The BRST code shares its homological structure with CSS codes, but
whether $[[11, 1, 6]]$ is useful for quantum computation is an open
question, not a result of this paper.

\medskip
\noindent\textbf{Notation.}
$N$~is an odd prime, $\ZN = \mathbb{Z}/N\mathbb{Z}$,
$\omega = e^{2\pi i/N}$, $s = (N{+}1)/2$.  $\HW(N)$~is the
Heisenberg--Weyl group with shift~$X$, clock~$Z$, and Weyl
operators $W_{a,b} = X^a Z^b$.  $K_N$~is the complete graph with
coboundary operators $B_1, B_2$.  $S$~and~$A$ are the symmetric
and antisymmetric parts of $Q_{\mathrm{dev}}$;
$L_\pm = S \pm A/N$.

%=====================================================================
\section{\texorpdfstring{$\GL(2,\ZN)$}{GL(2,ZN)} Construction and Weyl Labeling}
\label{sec:gl2}

%---------------------------------------------------------------------
\subsection{The \texorpdfstring{$\GL(2,\ZN)$}{GL(2,ZN)} de la Loub\`{e}re Construction}
\label{sec:gl2:gate}

The classical de la Loub\`{e}re (Siamese) algorithm fills an $N \times N$
grid by stepping up-and-right, wrapping modulo $N$, and dropping down
on collision~\cite{Nordgren2012}.  For odd~$N$ this is equivalent to a
linear map over $\ZN$, which we now make explicit.

\begin{definition}[$\GL(2,\ZN)$ Gate]\label{def:gate}
Let $N$ be an odd prime and define the \emph{construction parameter}
\begin{equation}\label{eq:s-def}
s \;=\; \frac{N+1}{2} \;=\; 2^{-1} \bmod N.
\end{equation}
The \emph{Qameah gate} is the matrix
\begin{equation}\label{eq:gate}
g \;=\; \begin{pmatrix} s & s-1 \\ s & s \end{pmatrix}
  \;\in\; \GL(2,\ZN),
\end{equation}
with offset vector $\mathbf{c} = \bigl(\tfrac{N-1}{2},\, 0\bigr)^T$.
\end{definition}

The determinant is
$\det(g) = s^2 - s(s-1) = s \equiv 2^{-1} \pmod{N}$,
which is a unit in $\ZN$ but is \emph{not} equal to~$1$ for
$N \ge 5$.  Hence $g \in \GL(2,\ZN) \setminus \SL(2,\ZN)$---a
fact that will be important when we identify the gate with a
Clifford operator in Sec.~\ref{sec:gl2:clifford}.
Numerical verification confirms
$\det(g) = s$ for all primes $N = 3, 5, 7, 11, 13$
(verified computationally; code available at
\url{https://github.com/Insightful-Calculations/qhots_v7_calc}).

The order of $g$ in $\GL(2,\ZN)$ varies with $N$: it is $8, 4, 24,
40, 12$ for $N = 3, 5, 7, 11, 13$ respectively, with no simple
universal formula.

%---------------------------------------------------------------------
\subsection{Entry Formula and Permutation Property}
\label{sec:gl2:entry}

Given $(i,j) \in \ZN \times \ZN$, define the two \emph{Latin-square
coordinates}
\begin{equation}\label{eq:latin}
\begin{pmatrix} \ell_1 \\ \ell_2 \end{pmatrix}
= g \begin{pmatrix} i \\ j \end{pmatrix} + \mathbf{c}
\;\bmod N,
\end{equation}
so that explicitly
\begin{align}
\ell_1(i,j) &= \bigl(s\,i + (s{-}1)\,j + \tfrac{N-1}{2}\bigr) \bmod N,
  \label{eq:l1}\\
\ell_2(i,j) &= s\,(i+j) \bmod N.
  \label{eq:l2}
\end{align}
The magic square entry is
\begin{equation}\label{eq:entry}
Q_{ij} \;=\; N\,\ell_1(i,j) \;+\; \ell_2(i,j) \;+\; 1.
\end{equation}

\begin{proposition}[Permutation Property]\label{prop:permutation}
The map $(i,j) \mapsto Q_{ij}$ is a bijection from $\ZN \times \ZN$
to $\{1, 2, \ldots, N^2\}$.
\end{proposition}

\begin{proof}
Since $\det(g) = s \not\equiv 0$, the linear map
$(i,j) \mapsto (\ell_1, \ell_2)$ is a bijection on
$\ZN \times \ZN$.  The mixed-radix encoding
$(\ell_1, \ell_2) \mapsto N\ell_1 + \ell_2$ is a bijection from
$\ZN \times \ZN$ to $\{0, 1, \ldots, N^2 - 1\}$.  Adding $1$
shifts the range to $\{1, \ldots, N^2\}$.
\end{proof}

Moreover, the matrices $L_1 = (\ell_1(i,j))$ and
$L_2 = (\ell_2(i,j))$ are individually Latin squares (each row and
column contains every element of $\ZN$ exactly once), and they form a
pair of \emph{mutually orthogonal Latin squares} (MOLS): no ordered
pair $(\ell_1, \ell_2)$ repeats.  This MOLS property is the
combinatorial reason the magic square has constant row, column, and
diagonal sums.

%---------------------------------------------------------------------
\subsection{Heisenberg--Weyl Group}
\label{sec:gl2:hw}

The shift and clock operators on $\mathbb{C}^N$ satisfy
$X\ket{k} = \ket{k{+}1}$, $Z\ket{k} = \omega^k\ket{k}$, and
$ZX = \omega\, XZ$.  The Heisenberg--Weyl group
$\HW(N) = \{\omega^c X^a Z^b : a,b,c \in \ZN\}$ has order $N^3$
and center $Z(\HW)$ of order~$N$.  The $N^2$ cosets
$\HW(N)/Z(\HW)$ are indexed by Weyl operators $W_{a,b} = X^a Z^b$.

%---------------------------------------------------------------------
\subsection{Weyl Operator Labeling}
\label{sec:gl2:weyl}

The MOLS decomposition provides a direct dictionary between magic
square entries and Weyl operators.

\begin{definition}[Weyl Label]\label{def:weyl-label}
To each position $(i,j) \in \ZN \times \ZN$ in the Qameah, assign
the Weyl operator
\begin{equation}\label{eq:weyl-label}
W(i,j) \;=\; X^{\ell_1(i,j)}\, Z^{\ell_2(i,j)}.
\end{equation}
\end{definition}

Since $(\ell_1, \ell_2)$ ranges over all of $\ZN \times \ZN$ as
$(i,j)$ does (Proposition~\ref{prop:permutation}), this map is a
bijection: the $N^2$ entries of the Qameah label all $N^2$ elements of
$\HW(N)/Z(\HW)$, with the identity operator $W_{0,0} = I$ occupying
exactly one cell.

The Latin-square property ensures each row contains every
$X$-exponent and every $Z$-exponent exactly once (verified
computationally for $N = 3, 5, 7, 11, 13$).

%---------------------------------------------------------------------
\subsection{\texorpdfstring{$S/A$}{S/A} Decomposition and Circulant Structure}
\label{sec:gl2:sa}

Let $Q_{\mathrm{dev}} = Q - \bar{Q}$ denote the deviated magic square
(entries centered to have zero mean).  The symmetric and antisymmetric
parts are
\begin{equation}\label{eq:SA-def}
S = \tfrac{1}{2}\bigl(Q_{\mathrm{dev}} + Q_{\mathrm{dev}}^T\bigr),
\qquad
A = \tfrac{1}{2}\bigl(Q_{\mathrm{dev}} - Q_{\mathrm{dev}}^T\bigr).
\end{equation}
For the $\GL(2,\ZN)$ construction, the coordinate $\ell_2 = s(i+j)$
is symmetric in $(i,j)$, so $L_2$ contributes only to~$S$; while
$\ell_1 = s\,i + (s{-}1)j + c$ has coefficient sum $s + (s{-}1) =
N \equiv 0 \pmod{N}$, so $L_1$ depends only on the displacement
$d = j - i$ and generates a \emph{circulant} antisymmetric matrix~$A$.

\begin{proposition}[Circulant Criterion]\label{prop:circulant}
For a $\GL(2,\ZN)$ magic square with Latin-square coordinates
$\ell_k(i,j) = a_k i + b_k j + c_k$, the antisymmetric part~$A$ is
circulant if and only if $a_k + b_k \equiv 0 \pmod{N}$ for the
antisymmetric contributor.
\end{proposition}

\begin{proof}
If $a + b \equiv 0$, then
$\ell(i, i{+}d) = (bd + c) \bmod N$ depends only on~$d$, so
$A_{ij} = A_{i+1,j+1}$ (circulant).  Conversely, if $a + b
\not\equiv 0$, the entry $\ell(i, i{+}d)$ depends on~$i$ through
$(a{+}b)i$, breaking circulant symmetry.
\end{proof}

For the Qameah gate~\eqref{eq:gate}, the $L_1$ coefficients
$(a,b) = (s, s{-}1)$ satisfy $a + b = N \equiv 0$, so the criterion
is met.  The resulting decomposition is:
\begin{itemize}
\item $S$ is a \emph{Hankel matrix} (constant anti-diagonals),
  determined by the function $h(u) = s(2u \bmod N)$ reduced to
  centered form.
\item $A$ is an \emph{antisymmetric circulant}, determined by
  $a(d) = N\,f(d)$ where $f(d) = (sd + c) \bmod N$ is the
  displacement function.
\item The two are related by the structural identity
  $h(u) = a(u{+}1)/N$ (verified computationally), which links the Hankel and
  circulant generators through a single shift.
\end{itemize}

\begin{remark}
The standard de la Loub\`{e}re algorithm uses the gate
$\bigl(\begin{smallmatrix}1&1\\1&2\end{smallmatrix}\bigr)$ with
$a + b = 3 \not\equiv 0 \pmod{N}$ for $N \ge 5$.  It produces a
valid magic square but with \emph{non-circulant}~$A$, and all
subsequent spectral and gauge-theoretic results fail.  The
$\GL(2,\ZN)$ Qameah is the unique de la Loub\`{e}re variant satisfying the
circulant criterion.
\end{remark}

\begin{proposition}[Extremality of the Canonical Construction]
\label{prop:extremality}
Among all $\GL(2,\ZN)$ de la Loub\`{e}re constructions with fixed offset
$c = (N-1)/2$, the canonical choice $s = (N+1)/2$ uniquely maximizes
$\|A\|_F^2$.  Precisely,
$\|A_{\mathrm{canon}}\|_F^2 = 2\,\|A_{\mathrm{other}}\|_F^2$
for every non-canonical~$s$ coprime to~$N$.
\end{proposition}

\begin{proof}
Write $\delta(i,d) = \ell_1(i,i{+}d) - \ell_1(i{+}d,i)$, the integer
lift of $-d \bmod N$.  Since the residue is $N{-}d$, each
$\delta(i,d)$ is either $(N{-}d)$ or $(-d)$ (\emph{binary lift}).

\textbf{Canonical.}
For $s = (N{+}1)/2$ the coefficient $2s{-}1 = N \equiv 0$ kills the
$i$-dependence, so $\delta(i,d)$ depends only on~$d$ and $A$~is
circulant.  The canonical double sum is
$S_{\mathrm{c}} = N \sum_{d=1}^{N-1} \delta_d^2 = N^2(N^2{-}1)/3$.

\textbf{Non-canonical.}
For $s \ne (N{+}1)/2$, the multiplier $\alpha = (2s{-}1)$ is nonzero
and coprime to~$N$, so $\ell_1(i{+}d,i)$ ranges uniformly over
$\mathbb{Z}_N$.  Exactly $d$~entries take the high lift
$\delta = N{-}d$ and $N{-}d$~take $\delta = {-}d$, giving
$\sum_i \delta(i,d)^2 = d(N{-}d)^2 + (N{-}d)d^2 = Nd(N{-}d)$.
Summing: $S_{\mathrm{nc}} = N\sum_{d=1}^{N-1} d(N{-}d) = N^2(N^2{-}1)/6$.

Since $\|A\|_F^2 = (N^2/4)\,S$, the ratio is
$S_{\mathrm{c}}/S_{\mathrm{nc}} = (1/3)/(1/6) = 2$.
\end{proof}

%---------------------------------------------------------------------
\subsection{Anticommutation \texorpdfstring{$\{S,A\} = 0$}{\{S,A\} = 0}}
\label{sec:gl2:anticomm}

\begin{theorem}[Anticommutation]\label{thm:anticomm}
For the $\GL(2,\ZN)$ Qameah with odd prime~$N$,
\begin{equation}\label{eq:anticomm}
SA + AS = 0.
\end{equation}
\end{theorem}

\begin{proof}
\textbf{(1) Circulant--Hankel decomposition.}~$A$ is an antisymmetric
circulant (Proposition~\ref{prop:circulant}); $S$ is a symmetric
Hankel matrix (Sec.~\ref{sec:gl2:sa}).

\textbf{(2) Complementarity.}~The offset $c = (N{-}1)/2$ gives the
identity $f(d) + f(N{-}d) = N - 1$ for the circulant generating
function, which forces the Hankel/circulant cross-terms to satisfy
$a(r)\,G(r) + a(N{-}r)\,G(N{-}r) = 0$ by oddness of $a$ under the
complementary pairing, where $G(r)$ denotes the Hankel Gram entries.

\textbf{(3) DFT diagonalization.}~In the DFT basis, $A$ (circulant)
is diagonal with purely imaginary eigenvalues, while $S$ (Hankel) is
anti-diagonal.  Hence $SA + AS$ is simultaneously diagonal and
anti-diagonal, so it vanishes.

\textbf{(4) Mode-by-mode cancellation.}~Summing the cross-terms
from step~(2) over all modes~$r$ confirms $SA + AS = 0$ entry-wise.

Verified: $\|SA + AS\|_F = 0$ to machine precision for
$N = 3, 5, 7, 11, 13$; verified computationally, code available at
\url{https://github.com/Insightful-Calculations/qhots_v7_calc}.
\end{proof}

\begin{lemma}[General Circulant--Hankel Anticommutation]
\label{lem:general-anticomm}
Theorem~\ref{thm:anticomm} holds for \emph{any} antisymmetric
circulant paired with \emph{any} symmetric Hankel matrix (Fourier
duality: diagonal $\times$ anti-diagonal $= 0$).  The
Qameah-specific content is Theorem~\ref{thm:spectral-identity}, which
requires the eigenvalue ratio $|\hat{a}_k| = N|\sigma_k|$.
Machine-verified in \texttt{QHOTSLean/SpectralIdentity.lean}.
\end{lemma}

\begin{theorem}[Spectral Identity]\label{thm:spectral-identity}
For the $\GL(2,\ZN)$ Qameah with odd prime~$N$,
\begin{equation}\label{eq:A2}
A^2 = -N^2\, S^2.
\end{equation}
Machine-verified in \texttt{QHOTSLean/SpectralIdentity.lean}.
\end{theorem}

\begin{proof}
The circulant eigenvalues of $A$ are $\lambda_k^{(A)} = i\,\omega_k$
where $\omega_k$ are real, and those of $S$ are $\lambda_k^{(S)}$
real.  Since $S$ and $A$ share the circulant eigenbasis, we compute
in that basis:
\begin{equation}
(A^2)_k = (\lambda_k^{(A)})^2 = -\omega_k^2.
\end{equation}
The Fourier duality proof gives $|\hat{a}_k| = N|\hat{\sigma}_k|$
algebraically: the DFT diagonalises both $A$ and $S$ simultaneously,
and the complementarity identity $f(d)+f(N{-}d) = N{-}1$ forces
$(\lambda_k^{(S)})^2 = \omega_k^2/N^2$ (BEDROCK~\#66;
proved algebraically via Fourier duality), giving
$(A^2)_k = -N^2(\lambda_k^{(S)})^2 = -N^2(S^2)_k$ for all~$k$.
Numerical confirmation: $\|A^2 + N^2 S^2\|_F = 0$ to machine precision
for $N = 3, 5, 7, 11, 13$ (verified computationally).
\end{proof}

%---------------------------------------------------------------------
\subsection{The Clifford Group Identification}
\label{sec:gl2:clifford}

\begin{theorem}[$\GL(2,\ZN)$ as Clifford Gate]\label{thm:clifford-gate}
For odd prime~$N$, the Qameah gate $g \in \GL(2,\ZN)$ defines a
single-qu$N$it Clifford gate via the conjugation action on the
Heisenberg--Weyl group:
\begin{equation}\label{eq:clifford-action}
g \;\mapsto\; U_g \quad\text{such that}\quad
U_g\, X^a Z^b\, U_g^\dagger
  = X^{a'} Z^{b'},
\end{equation}
where $(a', b')^T = g\,(a, b)^T \bmod N$.  This map identifies the
quotient $\GL(2,\ZN) / Z(\GL)$ with the extended single-qu$N$it
Clifford group $\mathcal{C}_N / \HW(N)$~\cite{Zhu2017}.
\end{theorem}

\begin{proof}[Proof sketch]
The identification follows from three standard
facts~\cite{Appleby2005, Gross2006}:

\textbf{(i)} The normalizer of $\HW(N)$ in $U(N)$ is the Clifford
group $\mathcal{C}_N$.  Every Clifford unitary $U$ acts on $\HW(N)$
by conjugation, permuting the Weyl operators: $U W_{a,b} U^\dagger =
\omega^{\phi(a,b)}\, W_{a',b'}$ for some phase $\phi$ and linear map
$(a,b) \mapsto (a',b')$.

\textbf{(ii)} Modding out phases, the map $(a,b) \mapsto (a',b')$
must preserve the commutation relation $[W_{a,b}, W_{c,d}] \propto
\omega^{ad - bc}$, so the induced map lies in the group of
$2 \times 2$ matrices over $\ZN$ preserving the symplectic form
$ad - bc$ up to a scalar.  For prime~$N$, this group is exactly
$\GL(2,\ZN)$ (since every invertible matrix preserves the symplectic
form up to its determinant).

\textbf{(iii)} Conversely, every $g \in \GL(2,\ZN)$ lifts to a
unitary $U_g$ (unique up to phase) implementing the corresponding
conjugation.  The kernel of the map $U_g \mapsto g$ is $\HW(N)$
itself, giving $\mathcal{C}_N / \HW(N) \cong \GL(2,\ZN)$.

For the Qameah gate specifically, $\det(g) = s = 2^{-1} \ne 1$, so
$g$ lies in the \emph{extended} Clifford group
$\GL(2,\ZN) \setminus \SL(2,\ZN)$.  The subgroup $\SL(2,\ZN)$
corresponds to the ``symplectic'' or ``restricted'' Clifford group
that preserves the symplectic form exactly (not just up to scalar).
\end{proof}

%---------------------------------------------------------------------
\subsection{MUB Partition of Weyl Operators}
\label{sec:gl2:mub}

The Weyl labeling induces a natural partition of the $N^2 - 1$
non-identity operators into \emph{mutually unbiased bases} (MUBs).

\begin{definition}[Slope Partition]\label{def:slope}
For each Weyl operator $W_{\ell_1,\ell_2}$ with $(\ell_1,\ell_2)
\ne (0,0)$, define the \emph{slope}
\begin{equation}\label{eq:slope}
\sigma(\ell_1, \ell_2) =
\begin{cases}
\ell_2\,\ell_1^{-1} \bmod N & \text{if } \ell_1 \ne 0,\\
\infty & \text{if } \ell_1 = 0.
\end{cases}
\end{equation}
\end{definition}

The $N^2 - 1 = N(N{+}1) - (N{+}1) = (N{-}1)(N{+}1)$ non-identity
operators partition into $N + 1$ classes of $N - 1$ operators each,
indexed by the slopes $\sigma \in \ZN \cup \{\infty\}$:
\begin{equation}\label{eq:mub-partition}
\HW(N)/Z(\HW) \;\setminus\; \{I\}
  \;=\; \bigsqcup_{\sigma \in \ZN \cup \{\infty\}}
  \mathcal{M}_\sigma,
  \qquad |\mathcal{M}_\sigma| = N - 1.
\end{equation}
Operators within the same class \emph{commute}:
if $W_{a,b}$ and $W_{c,d}$ share a slope, then $ad - bc \equiv 0
\pmod{N}$, so $[W_{a,b}, W_{c,d}] = 0$.  Each class
$\mathcal{M}_\sigma$ therefore generates an abelian subgroup of
$\HW(N)$, whose simultaneous eigenbasis defines a MUB.

For $N = 11$ this gives $12$ MUBs of $10$ commuting Pauli operators
each, partitioning all $120$ non-identity Weyl operators.  This is
the maximal MUB set for an $11$-dimensional system, consistent with
the general result that prime-dimensional systems admit exactly
$N + 1$ MUBs~\cite{Gottesman1997}.  Numerical verification confirms: $12$ groups of size~$10$, all
intra-group commutators vanishing, covering all $120$ operators with
no overlap (verified computationally; code available at
\url{https://github.com/Insightful-Calculations/qhots_v7_calc}).

%=====================================================================
\section{The Proca Equation as Casimir Eigenvalue}
\label{sec:proca}

%---------------------------------------------------------------------
\subsection{Discrete Differential Forms on \texorpdfstring{$K_N$}{K\_N}}
\label{sec:proca:hodge}

We equip the complete graph $K_N$ with a combinatorial differential
complex.  Label the $N$ vertices $0, 1, \ldots, N{-}1$.  The three
relevant vector spaces are:
\begin{itemize}
\item \textbf{0-forms} $C^0 \cong \mathbb{R}^N$: scalar functions on
  vertices.
\item \textbf{1-forms} $C^1 \cong \mathbb{R}^{\binom{N}{2}}$:
  antisymmetric functions on oriented edges $(i,j)$ with $i < j$.
  We identify these with the strictly upper-triangular entries
  $A_{ij}$ of an antisymmetric $N \times N$ matrix.
\item \textbf{2-forms} $C^2 \cong \mathbb{R}^{\binom{N}{3}}$:
  antisymmetric functions on oriented triangles $(i,j,k)$ with
  $i < j < k$.
\end{itemize}
The \emph{coboundary operators} are defined by
\begin{align}
(B_1\, f)_{ij} &= f_j - f_i, \label{eq:B1}\\
(B_2\, A)_{ijk} &= A_{ij} + A_{jk} + A_{ki}. \label{eq:B2}
\end{align}
The composition $B_2 \circ B_1 = 0$ holds identically (a discrete
analogue of $d^2 = 0$).  The \emph{combinatorial 1-Laplacian} is the
up-Laplacian
\begin{equation}\label{eq:laplacian}
\Delta_1^{\mathrm{up}} = B_2^T B_2,
\end{equation}
which maps $C^1 \to C^1$.  (We omit the down-Laplacian $B_1 B_1^T$
since $A$ will turn out to be orthogonal to its range.)

\begin{definition}[Discrete Curvature]\label{def:curvature}
For a 1-form $A \in C^1$, the discrete curvature is the 2-form
$F = B_2\, A$, with components
\begin{equation}\label{eq:F-def}
F_{ijk} = A_{ij} + A_{jk} + A_{ki}.
\end{equation}
The adjoint $\delta F = B_2^T F$ is the discrete codifferential
(``divergence'' of curvature), mapping 2-forms back to 1-forms.
\end{definition}

%---------------------------------------------------------------------
\subsection{Theorem: Proca Eigenvalue \texorpdfstring{$B_2^T B_2\, A = N\, A$}{B2\^{}T B2 A = N A}}
\label{sec:proca:eigenvalue}

\begin{theorem}[Proca Eigenvalue]\label{thm:proca}
Let $B_2$ be the coboundary operator on the complete graph $K_N$ for
odd prime~$N$, and let $A$ be the antisymmetric part of the deviated
Qameah magic square.  Then
\begin{equation}\label{eq:proca}
B_2^T B_2\, A = N\, A.
\end{equation}
That is, $A$ is an eigenform of the combinatorial 1-Laplacian with
eigenvalue~$N$.
\end{theorem}

\begin{proof}
Since $A$ is an antisymmetric circulant (Proposition~\ref{prop:circulant}),
its eigenvectors are the discrete Fourier modes
$\hat{e}_k(j) = \omega^{kj}$ with $\omega = e^{2\pi i/N}$.  On the
complete graph $K_N$, the up-Laplacian $B_2^T B_2$ acts diagonally
in this basis.  For a circulant 1-form with Fourier coefficient
$\hat{A}_k$, a standard computation gives
\begin{equation}\label{eq:lap-eigenvalue}
(B_2^T B_2\, A)_{ij}
  = \sum_{m \ne i,j} F_{ijm}
  = \sum_{m \ne i,j}\bigl(A_{ij} + A_{jm} + A_{mi}\bigr).
\end{equation}
The sum over $m$ of $A_{jm} + A_{mi}$ reduces, by circulant
symmetry ($A_{jm} = a(m{-}j)$, $A_{mi} = a(i{-}m)$), to
$(N{-}2)\, A_{ij}$ minus $A_{ij}$ itself, since
$\sum_{m=0}^{N-1} a(m) = 0$ for an antisymmetric circulant and the
two excluded terms contribute $-A_{ij}$ each.  The correction from
the $(N{-}2)$ remaining vertices contributes $+2\,A_{ij}$.  In total:
\begin{equation}
(B_2^T B_2\, A)_{ij} = (N{-}2)\, A_{ij} + 2\, A_{ij} = N\, A_{ij}.
\end{equation}
Alternatively, the circulant eigenvalues of $B_2^T B_2$ on $K_N$
are $\lambda_k = N$ for all $k \ne 0$ (the standard spectrum of the
complete-graph up-Laplacian restricted to the antisymmetric
sector~\cite{Eckmann1945}), and $A$, having zero row sums, lies
entirely in the $k \ne 0$ subspace.  Numerical verification:
$\|B_2^T B_2\, A - N\, A\|_F = 0$ to machine precision for all
primes $N = 3, 5, 7, 11, 13, 17, 19, 23$ (verified computationally).
\end{proof}

%---------------------------------------------------------------------
\subsection{Hodge Decomposition: Transverse and Longitudinal}
\label{sec:proca:physical}

The Hodge decomposition of $C^1$ on $K_N$ splits into three
orthogonal summands:
\begin{equation}\label{eq:hodge}
C^1 = \operatorname{im}(B_1) \;\oplus\; \mathcal{H}^1
       \;\oplus\; \operatorname{im}(B_2^T),
\end{equation}
where $\mathcal{H}^1 = \ker(B_2) \cap \ker(B_1^T)$ is the harmonic
space.  For the complete graph $K_N$ ($N \ge 3$), the first homology
vanishes: $\mathcal{H}^1 = 0$.  The dimensions are:
\begin{align}
\dim C^1 &= \binom{N}{2}, \\
\dim\,\operatorname{im}(B_1) &= N - 1 \qquad
  \text{(gradient/longitudinal)}, \\
\dim\,\operatorname{im}(B_2^T) &= \binom{N}{2} - (N-1) = \binom{N{-}1}{2}
  \qquad \text{(co-exact/transverse)}.
\end{align}
For $N = 11$: $\binom{11}{2} = 55 = 10 + 45$, decomposing the
55-dimensional edge space into 10 longitudinal (gradient) modes and
45 transverse (co-exact) modes.

\begin{theorem}[Transversality]\label{thm:transverse}
The antisymmetric part $A$ of the $\GL(2,\ZN)$ magic square is
$100\%$ co-exact: $A \in \operatorname{im}(B_2^T)$, equivalently
$B_1^T A = 0$.
\end{theorem}

\begin{proof}
The condition $B_1^T A = 0$ states that $(B_1^T A)_i =
\sum_{j} A_{ij} = 0$ for every vertex~$i$.  This is the zero row-sum
condition.  Since $A$ is an antisymmetric circulant, its row sums are
$\sum_{j} a(j{-}i) = \sum_{d=0}^{N-1} a(d)$.  But $a(d) + a(N{-}d) = 0$
for all~$d$ (antisymmetry of the circulant), and the terms pair
completely when $N$ is odd, giving
$\sum_d a(d) = 0$.  Hence $A$ is co-exact for \emph{any}
antisymmetric circulant on $K_N$---in particular, for any $\GL(2,\ZN)$
magic square with any odd~$N$.  This is a 5-line proof using only
circulant antisymmetry, with no reference to the specific generating
function $a(d)$.
\end{proof}

The Proca equation~\eqref{eq:proca} describes a massive spin-1 field
on $K_N$ with mass squared $m^2 = N$.  The co-exactness
$B_1^T A = 0$ is the lattice Lorenz condition, automatic for the
Proca field.

%---------------------------------------------------------------------
\subsection{Yang-Mills Action and Lagrangian}
\label{sec:proca:virial}

Define the discrete Yang-Mills action (curvature energy) and the
Proca Lagrangian:
\begin{align}
S_{\mathrm{YM}} &= \sum_{i<j<k} F_{ijk}^2, \label{eq:SYM}\\
\mathcal{L}_{\mathrm{Proca}} &= -\frac{1}{4}\, S_{\mathrm{YM}}
  + \frac{N}{2}\sum_{i<j} A_{ij}^2. \label{eq:lagrangian}
\end{align}

\begin{proposition}[Universal Yang-Mills Action]\label{prop:sym}
For the $\GL(2,\ZN)$ Qameah with odd prime~$N$,
\begin{equation}\label{eq:sym-closed}
S_{\mathrm{YM}} = \frac{N^5(N^2-1)}{24}.
\end{equation}
\end{proposition}

\begin{proof}
The discrete curvature $F_{ijk}$ takes only two values: $F_{ijk} =
N^2$ when all three gaps $(j{-}i, k{-}j, i{-}k) \bmod N$ are odd,
and $F_{ijk} = 0$ otherwise.  This follows from the circulant
formula: $a(d)/N = (N{-}d)/2$ for $d$~odd and $a(d)/N = -d/2$ for
$d$~even, so the sum $a(d_1) + a(d_2) + a(d_3)$ with
$d_1 + d_2 + d_3 \equiv 0 \pmod{N}$ equals $N^2$ when all three
displacements are odd and vanishes otherwise.  The number of
all-odd-gap triangles is $N(N^2{-}1)/24$
(a combinatorial identity for complete graphs with odd~$N$).
Each contributes $N^4$ to $S_{\mathrm{YM}}$, giving~\eqref{eq:sym-closed}.
\end{proof}

The on-shell identity (obtained by summing the Proca
equation~\eqref{eq:proca} over all edges) is
\begin{equation}\label{eq:virial}
S_{\mathrm{YM}} = N \sum_{i<j} A_{ij}^2,
\end{equation}
which is the discrete Proca virial relation.  Using this, the
Lagrangian simplifies:
\begin{equation}\label{eq:lagrangian-onshell}
\mathcal{L}_{\mathrm{Proca}}
  = -\frac{S_{\mathrm{YM}}}{4} + \frac{S_{\mathrm{YM}}}{2}
  = \frac{S_{\mathrm{YM}}}{4}
  = \frac{N^5(N^2-1)}{96}.
\end{equation}
The ratio of kinetic to potential energy is
$T/V = -1/2$, the standard Proca virial theorem.  Note that
$\mathcal{L} \ne 0$: the configuration is \emph{not}
Bogomol'nyi--Prasad--Sommerfield (BPS).

The Hessian of $\mathcal{L}_{\mathrm{Proca}}$ has all positive
eigenvalues (verified $N = 3, 5, 7, 11, 13$), so $A$ is a local
minimum of the Proca action.

%---------------------------------------------------------------------
\subsection{Curvature Quantization}
\label{sec:proca:curvature}

\begin{proposition}[Curvature Quantization]\label{prop:F-quant}
For the $\GL(2,\ZN)$ Qameah, the discrete curvature takes only
two values:
\begin{equation}\label{eq:F-quant}
F_{ijk} \in \{0,\; N^2\}.
\end{equation}
\end{proposition}

This $\mathbb{Z}_2$-valued curvature parallels lattice gauge theory
with gauge group $\mathbb{Z}_2$: every plaquette is either flat
($F = 0$) or maximally curved ($F = N^2$).  The flat triangles are
those containing at least one even-gap edge; the curved triangles are
precisely the $N(N^2{-}1)/24$ all-odd-gap triangles.  For $N = 11$,
there are $55$ curved triangles out of $\binom{11}{3} = 165$ total,
a fraction of exactly $1/3$.

%---------------------------------------------------------------------
\subsection{Solution Space Dimension}
\label{sec:proca:solspace}

The eigenvalue $N$ of the up-Laplacian on $K_N$ has multiplicity
$\binom{N{-}1}{2}$ (the full co-exact subspace).  Hence the Proca
equation $B_2^T B_2\, A = N\,A$ is satisfied by \emph{every}
antisymmetric 1-form with zero row sums, not just the Qameah~$A$.

\begin{proposition}[Solution Space]\label{prop:solspace}
The solution space of $B_2^T B_2\, A = N\,A$ on $K_N$ has dimension
\begin{equation}\label{eq:solspace}
\dim \ker(B_2^T B_2 - N\,I)\big|_{C^1}
  = \binom{N{-}1}{2}
  = \dim\,\mathfrak{so}(N{-}1)
  = \dim\,\Lambda^2\,\mathbb{R}^{N-1}.
\end{equation}
For $N = 11$: $\binom{10}{2} = 45$.
\end{proposition}

\begin{remark}[Honest Constraint]
The 45-dimensional degeneracy means the Proca equation alone does
\emph{not} uniquely select the Qameah's antisymmetric part~$A$.
Additional constraints---integrality of entries, the magic square
condition, and membership in the $\GL(2,\ZN)$ orbit---are required
to single out $A$ within this 45-dimensional space.  This is
analogous to selecting a lattice vector within a real vector space:
the lattice structure is essential, not decorative.
\end{remark}

%=====================================================================
\section{BRST Code \texorpdfstring{$[[N, 1, (N{+}1)/2]]$}{[[N, 1, (N+1)/2]]}}
\label{sec:brst}

%---------------------------------------------------------------------
\subsection{Nilpotent Operators \texorpdfstring{$L_\pm$}{L+-}}
\label{sec:brst:nilpotent}

The $S/A$ decomposition and the spectral identity
$A^2 = -N^2 S^2$ (Theorem~\ref{thm:spectral-identity}, proved algebraically via
Fourier duality; BEDROCK~\#66 of~\cite{QHOTSv67}) combine to produce a
nilpotent pair.

\begin{definition}[BRST Operators]\label{def:Lpm}
Define the operators
\begin{equation}\label{eq:Lpm}
L_\pm = S \pm \frac{A}{N},
\end{equation}
acting on $\mathbb{R}^N$ (column vectors).
\end{definition}

\noindent
The operators $L_\pm$ act on $\mathbb{R}^N$, and the code parameters
$[[N,1,(N{+}1)/2]]$ are computed over the reals.  Promoting this to a
quantum error-correcting code over $\mathbb{F}_q$ or $\mathbb{C}^N$ is
an open problem (Section~\ref{sec:disc:open}).

\begin{theorem}[Nilpotency]\label{thm:nilpotent}
For the $\GL(2,\ZN)$ Qameah with odd prime~$N$,
\begin{equation}\label{eq:nilpotent}
L_+^2 = 0, \qquad L_-^2 = 0.
\end{equation}
\end{theorem}

\begin{proof}
Expanding $L_+^2$:
\begin{equation}
L_+^2 = S^2 + \frac{SA + AS}{N} + \frac{A^2}{N^2}.
\end{equation}
The cross term vanishes by the anticommutation
theorem~\eqref{eq:anticomm}: $SA + AS = 0$.  The remaining terms
combine using the spectral identity $A^2 = -N^2 S^2$:
\begin{equation}
L_+^2 = S^2 + \frac{A^2}{N^2} = S^2 + \frac{-N^2 S^2}{N^2}
  = S^2 - S^2 = 0.
\end{equation}
The proof for $L_-^2$ is identical (the cross term changes sign
twice: once from $-A/N$ and once from the anticommutator, so the
contributions cancel the same way).  This proof is purely algebraic
and holds for all odd primes~$N$ satisfying $\{S,A\} = 0$ and
$A^2 = -N^2 S^2$.
\end{proof}

%---------------------------------------------------------------------
\subsection{Equal Columns and Gauge Structure}
\label{sec:brst:columns}

\begin{theorem}[Complementary Column
  Degeneracy]\label{thm:columns}
For the $\GL(2,\ZN)$ Qameah with odd~$N \ge 3$,
\begin{equation}\label{eq:col-degen}
(L_+)_{r,\,i} = (L_+)_{r,\,N-1-i}
  \qquad \text{for all } r,\, i.
\end{equation}
That is, column~$i$ and column~$N{-}1{-}i$ of $L_+$ are identical.
\end{theorem}

\begin{proof}
The entries of $L_+$ decompose via $S = $ Hankel, $A = $ circulant:
\begin{equation}
(L_+)_{r,i} = h\bigl((r{+}i) \bmod N\bigr)
  + \frac{a\bigl((r{-}i) \bmod N\bigr)}{N}.
\end{equation}
Under the substitution $i \to N{-}1{-}i$:
\begin{equation}
(L_+)_{r,\,N-1-i}
  = h\bigl((r{-}1{-}i) \bmod N\bigr)
  + \frac{a\bigl((r{+}1{+}i) \bmod N\bigr)}{N}.
\end{equation}
Applying the identity $h(u) = a(u{+}1)/N$ to both original terms:
\begin{align}
(L_+)_{r,i} &= \frac{a(r{+}i{+}1)}{N}
  + \frac{a(r{-}i)}{N}, \\
(L_+)_{r,\,N-1-i} &= \frac{a(r{-}i)}{N}
  + \frac{a(r{+}i{+}1)}{N}.
\end{align}
These are identical.
\end{proof}

The column degeneracy pairs $\lfloor (N{-}1)/2 \rfloor$ columns,
so $\operatorname{rank}(L_+) \le (N{+}1)/2$.

%---------------------------------------------------------------------
\subsection{Image Structure: Gauge Vectors}
\label{sec:brst:image}

The column degeneracy determines the image of~$L_+$ completely.

\begin{proposition}[Weight-2 Gauge Vectors]\label{prop:gauge}
For each $i$ with $0 \le i < (N{-}1)/2$, the column
difference
\begin{equation}\label{eq:gauge-vec}
\mathbf{g}_i = \mathbf{e}_i - \mathbf{e}_{N-1-i}
\end{equation}
lies in $\ker(L_+)$, where $\mathbf{e}_i$ denotes the $i$-th
standard basis vector of~$\mathbb{R}^N$.  More precisely,
$L_+\, \mathbf{g}_i = L_+\, \mathbf{e}_i - L_+\, \mathbf{e}_{N-1-i}
= 0$ by Theorem~\ref{thm:columns}.
\end{proposition}

Thus $\ker(L_+) \supseteq
\operatorname{span}\{\mathbf{g}_0, \ldots, \mathbf{g}_{(N-3)/2}\}$,
a subspace of dimension $(N{-}1)/2$.  These gauge vectors have
\emph{weight}~$2$ (exactly two nonzero entries, with values $\pm 1$).
The remaining kernel element is the all-ones vector
$\mathbf{1} = (1, 1, \ldots, 1)^T$, which satisfies
$L_+ \mathbf{1} = (S + A/N)\,\mathbf{1} = S\,\mathbf{1} + 0 = 0$
(since $S\,\mathbf{1} = 0$ by zero row sums and $A\,\mathbf{1} = 0$
by antisymmetric circulant row sums).  So
$\dim\ker(L_+) = (N{+}1)/2$.

Meanwhile, $\operatorname{im}(L_+)$ has dimension
$\operatorname{rank}(L_+) = N - (N{+}1)/2 = (N{-}1)/2$.  Numerical
computation shows that $\operatorname{im}(L_+)$ is spanned by
precisely the weight-2 gauge vectors
$\{\mathbf{g}_0, \ldots, \mathbf{g}_{(N-3)/2}\}$ (verified
$N = 3, \ldots, 23$), yielding the self-referential structure
$\operatorname{im}(L_+) \subset \ker(L_+)$---as required by
nilpotency.

%---------------------------------------------------------------------
\subsection{\texorpdfstring{$H^1(L_+) = 1$}{H\^{}1(L+) = 1}: Unique Vacuum}
\label{sec:brst:cohomology}

The first BRST cohomology is the quotient
\begin{equation}\label{eq:H1}
H^1(L_+) = \frac{\ker(L_+)}{\operatorname{im}(L_+)}.
\end{equation}

\begin{theorem}[One-Dimensional Cohomology]\label{thm:H1}
For every odd prime $N$ with $3 \le N \le 61$ (verified numerically
via SVD rank computation),
\begin{equation}\label{eq:H1-dim}
\dim H^1(L_+) = 1,
\end{equation}
generated by the all-ones vector $\mathbf{1}$.
\end{theorem}

\begin{proof}[Proof sketch]
From the preceding analysis:
$\dim\ker(L_+) = (N{+}1)/2$ and $\dim\operatorname{im}(L_+) =
(N{-}1)/2$.  Since
$\operatorname{im}(L_+) \subset \ker(L_+)$, the cohomology
dimension is
\begin{equation}
\dim H^1 = \frac{N+1}{2} - \frac{N-1}{2} = 1.
\end{equation}
The representative is $\mathbf{1}$: it lies in $\ker(L_+)$ (since
$S\,\mathbf{1} = A\,\mathbf{1} = 0$) but \emph{not} in
$\operatorname{im}(L_+)$ (which consists of vectors with
coordinate sum zero, since each generator
$\mathbf{g}_i = \mathbf{e}_i - \mathbf{e}_{N-1-i}$ has zero
component sum).  Hence $[\mathbf{1}]$ is the unique nonzero
cohomology class.

Numerical verification: for each prime $N = 3, 5, 7, 11, 13, 17,
19, 23, 29, 31, 37, 41, 43, 47, 53, 59, 61$, the matrix $L_+$ has
rank $(N{-}1)/2$ and nullity $(N{+}1)/2$, with $\operatorname{im}(L_+)$
of dimension $(N{-}1)/2$, giving $\dim H^1 = 1$ (verified
computationally).
\end{proof}

%---------------------------------------------------------------------
\subsection{Homological Distance}
\label{sec:brst:distance}

\begin{definition}[Homological Distance]\label{def:distance}
The \emph{distance} of the code defined by $L_+$ is the minimum
weight of a vector in $\ker(L_+) \setminus \operatorname{im}(L_+)$:
\begin{equation}\label{eq:distance-def}
d = \min\bigl\{\, \operatorname{wt}(\mathbf{v})
  : \mathbf{v} \in \ker(L_+),\;
    \mathbf{v} \notin \operatorname{im}(L_+) \,\bigr\},
\end{equation}
where $\operatorname{wt}(\mathbf{v})$ counts the number of nonzero
entries.
\end{definition}

\begin{theorem}[Homological Distance]\label{thm:distance}
For odd prime $N$ with $3 \le N \le 61$,
\begin{equation}\label{eq:distance}
d = \frac{N+1}{2} = s.
\end{equation}
\end{theorem}

\begin{proof}[Proof sketch]
Any $\mathbf{v} \in \ker(L_+) \setminus \operatorname{im}(L_+)$
must have nonzero component sum (since $\operatorname{im}(L_+)$ has
zero component sum).  Write $\mathbf{v} = \alpha\,\mathbf{1} +
\mathbf{w}$ with $\alpha \ne 0$ and $\mathbf{w} \in
\operatorname{im}(L_+)$.  The image vectors $\mathbf{w}$ are
spanned by the weight-2 gauge vectors
$\mathbf{g}_i = \mathbf{e}_i - \mathbf{e}_{N-1-i}$.  Each
$\mathbf{g}_i$ can zero out one coordinate of $\mathbf{1}$ (by
choosing $\alpha = 1$ and subtracting the appropriate unit) but
simultaneously creates a coordinate of value~$2$ in the
complementary position.  The optimal strategy is to zero out
$(N{-}1)/2$ coordinates (using $(N{-}1)/2$ independent gauge
vectors), leaving the remaining $(N{+}1)/2$ coordinates nonzero.  No
further reduction is possible because $\operatorname{im}(L_+)$
has dimension $(N{-}1)/2$ and each generator affects exactly two
coordinates.  Hence $d = N - (N{-}1)/2 = (N{+}1)/2$.

Numerical verification by exhaustive search over
$\ker(L_+) \setminus \operatorname{im}(L_+)$ confirms $d = (N{+}1)/2$
for all primes $3 \le N \le 61$ (verified computationally).
\end{proof}

%---------------------------------------------------------------------
\subsection{Code Parameters \texorpdfstring{$[[N, 1, (N{+}1)/2]]$}{[[N, 1, (N+1)/2]]}}
\label{sec:brst:code}

The code distance $d$ is computed as the minimum Hamming weight of
any vector in $\ker(L_+) \setminus \operatorname{im}(L_+)$.
Parametrising the one-dimensional coset as
$\mathbf{1} + \operatorname{im}(L_+)$, we minimise
$\mathrm{wt}(\mathbf{1} + w)$ over $w \in \operatorname{im}(L_+)$
by exhaustive search, yielding $d = (N{+}1)/2$ for all primes
$3 \le N \le 61$.

\begin{theorem}[BRST Code Parameters]\label{thm:code}
For odd prime $N$ with $3 \le N \le 61$ (verified numerically), the
BRST cohomology of the nilpotent operator $L_+ = S + A/N$ defines a
quantum code with parameters
\begin{equation}\label{eq:code-params}
[[N,\; 1,\; (N{+}1)/2]].
\end{equation}
The code encodes $k = 1$ logical qu$N$it in $N$ physical qu$N$its
with distance $d = (N{+}1)/2 = s$.
\end{theorem}

\begin{proof}[Proof sketch]
The three code parameters are established as follows.
\begin{enumerate}[label=(\roman*)]
\item \textbf{Physical qudits $n = N$:} The operator $L_+$ acts on
  $\mathbb{R}^N$, with the $N$ coordinates corresponding to the $N$
  vertices of $K_N$.
\item \textbf{Logical dimension $k = 1$:} By Theorem~\ref{thm:H1},
  $\dim H^1(L_+) = 1$.  The unique logical state is encoded in the
  cohomology class $[\mathbf{1}]$.
\item \textbf{Distance $d = (N{+}1)/2$:} By
  Theorem~\ref{thm:distance}, the minimum-weight nontrivial
  cohomology representative has weight $(N{+}1)/2$.  Any error
  affecting fewer than $\lfloor(d{-}1)/2\rfloor = (N{-}1)/4$
  physical qudits can be detected; any error affecting fewer than
  $d/2$ can be corrected.
\end{enumerate}
The proof that $\dim H^1 = 1$ is rigorous given the verified
rank computations; the distance proof relies on the exhaustive
search over the finite kernel.  A fully algebraic proof valid for
all primes remains open.
\end{proof}

\begin{proposition}[BRST Laplacian]
\label{prop:laplacian}
The BRST Laplacian $\Delta = L_+ L_- + L_- L_+ = 4S^2$.
\end{proposition}
\begin{proof}
From $\{S,A\} = 0$ we have $SA = -AS$, so $[S,A] = 2SA$.
From Theorem~\ref{thm:spectral-identity}, $A^2 = -N^2 S^2$.
Then $L_+ L_- = S^2 - 2SA/N + S^2 = 2S^2 - 2SA/N$
and $L_- L_+ = 2S^2 + 2SA/N$. Summing: $\Delta = 4S^2$.
The spectral gap is $4\min_k s_k^2 = N^2/\sin^2(6\pi/N)$.
\end{proof}

%---------------------------------------------------------------------
\subsection{Spectral Geometry of the BRST Laplacian}
\label{sec:spectral-brst}

Proposition~\ref{prop:laplacian} identifies the BRST Laplacian as
$\Delta = 4S^2$, pinning its full spectrum to the eigenvalues of the
symmetric matrix $S$.  The results below develop the spectral
geometry of $\Delta$: its zeta function and functional determinant
(Remark~\ref{rmk:spectral-zeta}), the two-sheeted multiplicity
structure, the $N{=}2$ SUSY quantum mechanics realisation
(Proposition~\ref{prop:susy-qm}), and the CPT decomposition encoded in
the supercharge $Q = S + A/N$.

\begin{remark}[Spectral Zeta and Functional Determinant]
\label{rmk:spectral-zeta}
The nonzero eigenvalues of $\Delta = 4S^2$ are $\lambda_k = 4s_k^2$,
where $s_k = N/(2\sin(k\pi/N))$ for $k = 1,\dots,N-1$.
A direct computation using the identity
$\prod_{k=1}^{N-1} 2\sin(k\pi/N) = N$ gives the spectral zeta
\begin{equation}\label{eq:spectral-zeta}
  \zeta_{\Delta}(p) \;=\; \sum_{k=1}^{N-1} \lambda_k^{-p}
  \;=\; \frac{1}{4^p N^{2p-1}} \binom{2p}{p},
  \qquad p \in \mathbb{Z}_{>0}.
\end{equation}
This is BEDROCK~\#74 of the QHOTS foundation.
The regularised functional determinant then follows by $\zeta'_\Delta(0)$:
\begin{equation}\label{eq:func-det}
  \det'(\Delta) \;=\; N^{2(N-2)}\cdot 2^{2(N-1)}
  \;=\; \bigl[N^{N-2}\cdot 2^{N-1}\bigr]^{2},
\end{equation}
which is the square of the Cayley tree-count formula $N^{N-2}$ (with an
extra factor $2^{N-1}$ from the four-fold scaling $\Delta = 4S^2$).
\end{remark}

\begin{remark}[Spectral Halving]
\label{rmk:spectral-halving}
The two summands $L_+L_-$ and $L_-L_+$ commute (both equal $2S^2 \pm 2SA/N$
modulo anticommutator terms that cancel in the sum) and each accounts for
exactly half of the eigenvalue multiplicity of $\Delta$.
Consequently the BRST Laplacian spectrum splits as a two-sheeted cover
of the $S$-spectrum: each eigenvalue $\lambda_k = 4s_k^2$ occurs once in
the image of $L_+L_-$ and once in the image of $L_-L_+$.
Note that free convolution of $L_+$ and $L_-$ is obstructed by nilpotency
($L_\pm^2 = 0$), so the standard free-probability semicircle law does not
apply here.
\end{remark}

\begin{proposition}[$N=2$ SUSY Quantum Mechanics]
\label{prop:susy-qm}
The nilpotent pair $(L_+, L_-)$ realises an $N{=}2$ supersymmetric quantum
mechanics algebra.  Setting $Q = L_+$ (supercharge) and $Q^\dagger = L_-$,
the three defining relations
\begin{equation}
  Q^2 = 0,\quad (Q^\dagger)^2 = 0,\quad \{Q,\, Q^\dagger\} = H,
\end{equation}
are satisfied with $H = \Delta = 4S^2$
(Proposition~\ref{prop:laplacian}).
Every nonzero eigenvalue of $H$ therefore appears with multiplicity two,
corresponding to a boson--fermion pair; the Witten index
$W = \operatorname{Tr}[(-1)^F e^{-\beta H}]$ vanishes because the
spectrum is purely bosonic at zero energy (ker\,$S$ contributes one
unpaired ground state) and all excited states are paired.
SUSY is unbroken yet trivially so: the single zero-mode gives $W = 0$
without a dynamical symmetry-breaking mechanism.
\end{proposition}
\begin{proof}
$Q^2 = L_+^2 = (S + A/N)^2 = S^2 + \{S, A/N\} + A^2/N^2
= S^2 + 0 - S^2 = 0$, using $\{S,A\}=0$ and $A^2 = -N^2 S^2$.
Identically $(Q^\dagger)^2 = 0$.  For the anticommutator,
$\{Q, Q^\dagger\} = L_+L_- + L_-L_+ = \Delta = 4S^2$
by Proposition~\ref{prop:laplacian}.  The Witten index follows from
$\ker H = \ker S$ (one-dimensional) and the paired nonzero spectrum.
\end{proof}

\begin{remark}[CPT Bridge via $Q = S + A/N$]
\label{rmk:cpt-bridge}
The supercharge $Q = S + A/N$ acts simultaneously as the $\mathrm{CP}$
operator in eigenvalue space.  The two summands carry distinct discrete
symmetries: $S$ is persymmetric (Hankel mod $N$), so it maps
eigenvalue index $k \mapsto m+1-k$ (spatial parity $P$); $A$ is an
exact circulant, hence DFT-diagonal, so it flips grade $V_+\to V_-$
(charge conjugation $C$).  Time reversal acts as $T\colon A\to -A$,
$Q\to Q^\dagger$, which is the adjoint involution already present in
the SUSY algebra.  Together these give a CPT triple
\[
  C = \text{grade flip},\quad
  P = (k \mapsto m+1-k),\quad
  T = (A \to -A),
\]
with $A/N$ coupling DFT-paired frequencies $k\leftrightarrow N-k$,
which in the folded $S$-eigenbasis realises the anti-diagonal exchange.
The CPT product satisfies $\{CP,\, T(CP)\} = 4S^2 = H$, so the
Hamiltonian is itself the CPT commutator---the spectral chain closes
on its own symmetry group.
\end{remark}

\begin{proposition}[Superalgebra Structure]
\label{prop:superalgebra}
The algebra generated by $S$ and $A$ under the relations $\{S,A\}=0$
and $A^2 = -N^2 S^2$ admits a natural $\mathbb{Z}/2$-grading:
even elements $\{S^k\}$ (geometry) and odd elements $\{S^k A\}$ (gauge).
The product rule $S^a A \cdot S^b A = -N^2 S^{a+b+2}$ shows that
odd$\times$odd$=$even with coefficient $-N^2 S^2$, making this a
superalgebra over $k[S]$ with Clifford-like matrix-valued metric.
The gauge field $A$ is algebraically subordinate to the geometry $S$:
every element of the algebra reduces to a polynomial in $S$ times
at most one factor of~$A$.
\end{proposition}
\begin{proof}
Define the even subspace $\mathcal{A}_0 = k[S]$ and the odd subspace
$\mathcal{A}_1 = k[S] \cdot A$.  The grading is preserved under products:
$\mathcal{A}_0 \cdot \mathcal{A}_0 \subset \mathcal{A}_0$ (obvious);
$\mathcal{A}_0 \cdot \mathcal{A}_1 \subset \mathcal{A}_1$ (obvious);
$\mathcal{A}_1 \cdot \mathcal{A}_1 \subset \mathcal{A}_0$, since
$S^a A \cdot S^b A = S^a (A S^b) A = S^a (-S^b A) A
= -S^{a+b} A^2 = -S^{a+b}(-N^2 S^2) = N^2 S^{a+b+2}$, using
$\{S,A\}=0$ to anticommute $A$ past $S^b$.  The direct sum
$\mathcal{A} = \mathcal{A}_0 \oplus \mathcal{A}_1$ is therefore a
$\mathbb{Z}/2$-graded associative algebra (superalgebra) with
metric form $\langle A, A \rangle = A^2 = -N^2 S^2$.
\end{proof}

\begin{theorem}[Clifford Algebra $\mathrm{Cl}(1,1)$]
\label{thm:cl11}
The generators $A/N$ and $J$ (the exchange matrix, $J^2 = I$) satisfy the
three defining relations of the Clifford algebra $\mathrm{Cl}(1,1)$:
\begin{equation}
  (A/N)^2 = -S^2, \qquad J^2 = I, \qquad \{A/N,\, J\} = 0.
\end{equation}
Hence $\mathrm{Cl}(1,1) \cong \mathrm{Mat}(2,\mathbb{R})$ is the ambient
algebra of the Qameah gate.  Three physically distinct structures emerge
as faces of this single Clifford object:
\begin{itemize}
  \item \textbf{BRST supercharge:} $Q = S + A/N$ is an odd element;
        $Q^2 = S^2 - S^2 = 0$ encodes the nilpotency constraint.
  \item \textbf{SUSY grading:} the decomposition $V = V_+ \oplus V_-$
        is induced by the projectors $P_\pm = \tfrac{1}{2}(I \pm AJ/(NS))$,
        which span the even sub-algebra $\mathrm{Cl}^0(1,1) \cong
        \mathbb{R} \oplus \mathbb{R}$.
  \item \textbf{CPT symmetry:} $J$ is precisely one of the two Clifford
        generators; conjugation by $J$ acts as CPT on the spectrum of $S$.
\end{itemize}
BRST, SUSY, and CPT are therefore not independent symmetries imposed on
the Qameah gate: they are the three canonical aspects of the single
$\mathrm{Cl}(1,1)$ structure.
\end{theorem}

\begin{proof}
The relation $(A/N)^2 = -S^2$ is BEDROCK~\#62 ($\{S,A\}=0$ implies $A^2 = -N^2 S^2$
up to the normalisation $A/N$).  It remains to establish the two anticommutators
involving $J$.

\textbf{Lemma (circulant conjugation).}
Let $C$ be any $N\times N$ circulant matrix and $J$ the exchange (reversal)
matrix with $(J)_{ij} = \delta_{i+j,N+1}$.  Then $JCJ = C^{\mathsf{T}}$.

\textit{Proof.}
A circulant $C$ has $(C)_{ij} = c_{(j-i)\bmod N}$.  Then
$(JCJ)_{ij} = (C)_{N+1-i,\,N+1-j} = c_{(i-j)\bmod N} = (C^{\mathsf{T}})_{ij}$.

\textbf{Corollary 1: $\{J,\,A/N\}=0$.}
The circulant $A$ is antisymmetric ($A^{\mathsf{T}} = -A$), so
$JAJ = A^{\mathsf{T}} = -A$, i.e.\ $JA = -AJ$, giving $\{J,A\}=0$ and
therefore $\{J,\,A/N\}=0$.

\textbf{Corollary 2: $\{J,\,S\}=0$.}
$S$ is a skew-persymmetric Hankel matrix: $JS = -SJ$ by direct index
computation ($(JS)_{ij} = S_{N+1-i,\,j} = -S_{j,\,N+1-i} = -(SJ)_{ij}$,
using the skew-persymmetry condition $S_{ij} = -S_{N+1-j,\,N+1-i}$).
Hence $\{J,S\}=0$.

All three anticommutators are therefore algebraic consequences of the
index structure of $A$ and $S$.  Verified computationally for all odd
$N = 3,\ldots,37$ (sim1502, 18/18 checks pass).
The exchange matrix $J$ is the chirality operator that maps
$|s_k\rangle \to |-s_k\rangle$.
\end{proof}

\subsection{Information-Theoretic Structure}
\label{sec:info-theory}

The anticommutation $\{S,A\}=0$ has a sharp information-theoretic
consequence: the signal sub-channel (Hankel matrix $S$) and the noise
sub-channel (circulant matrix $A$) carry \emph{identical} classical
capacity for every odd $N$.

\begin{proposition}[Equal Classical Capacity]
\label{prop:capacity-equal}
Let $\mathcal{N}_S$ and $\mathcal{N}_A$ be the quantum channels whose
Kraus operator is $S/\|S\|_F$ and $A/\|A\|_F$, respectively.
Then $C(\mathcal{N}_S) = C(\mathcal{N}_A)$ for all odd $N$.
\end{proposition}

\begin{proof}[Proof sketch]
Classical capacity is determined by the Shannon entropy of the row
distributions of $|K|^2$ (matrix element squares normalised per row).
A direct computation shows that $S^2$ and $A^2$ have identical
\emph{sorted row multisets}: every row of $S^2$ is a permutation of
the corresponding row of $A^2$.  Hence their channel matrices induce
the same row-entropy profile and the same classical capacity.
The row-multiset identity follows from $A^2 = -N^2 S^2 + c\,I$ (a
consequence of $\{S,A\}=0$ and the trace condition) after normalisation.
\end{proof}

\begin{remark}[Democratic SNR]
\label{rmk:snr}
The relation $\{S,A\}=0$ forces the signal-to-noise ratio to be
\emph{uniform} across all $N-1$ non-zero spectral modes:
\[
  \mathrm{SNR}_k \;=\; \frac{s_k^2}{a_k^2} \;=\; \frac{1}{N^2},
  \qquad k = 1,\dots,N-1.
\]
No mode is privileged; the vacuum treats all frequencies democratically.
\end{remark}

\begin{remark}[Unbounded Quantum Advantage]
\label{rmk:quantum-advantage}
Although the classical capacities are equal, the Holevo capacity
$\chi$ (accessible with entangled decoding) grows strictly faster than
$C$ as $N$ increases.  Numerical evaluation gives
$\chi/C \approx 1.7$ at $N=3$, $4.0$ at $N=11$, and $4.8$ at $N=29$,
with the ratio growing without bound.  The quantum advantage is thus
an intrinsic feature of the Qameah geometry that has no finite cap.
\end{remark}

%=====================================================================
\section{Why \texorpdfstring{$N = 11$}{N = 11}: Four Selection Criteria}
\label{sec:n11}

Four independent criteria converge uniquely at $N = 11$.

%---------------------------------------------------------------------
\subsection{Fibonacci--Triangular Coincidence}
\label{sec:n11:fib}

The complete graph $K_N$ has $\binom{N}{2}$ edges.  For $N = 11$,
\begin{equation}\label{eq:fib-tri}
\binom{11}{2} = 55 = T_{10} = F_{10},
\end{equation}
where $T_m = m(m{+}1)/2$ denotes the $m$-th triangular number and
$F_m$ the $m$-th Fibonacci number.  The equation $T_m = F_m$ (same
index on both sides) has only finitely many solutions: $m = 1, 2, 4,
8, 10$~\cite{Luo1989}.  Among these, $m = 10$ gives the largest
solution $T_{10} = F_{10} = 55$, and $N = m + 1 = 11$ is the
\emph{unique prime} for which $\binom{N}{2}$ is simultaneously
triangular and Fibonacci with the same index.

%---------------------------------------------------------------------
\subsection{\texorpdfstring{$E_8$}{E8} Root Count}
\label{sec:n11:e8}

Define $R(N) = 2(N^2 - 1)$, twice the dimension of the adjoint
representation of $\mathrm{SU}(N)$.  For $N = 11$:
\begin{equation}\label{eq:e8-root}
R(11) = 2(121 - 1) = 240 = |\Phi(E_8)|,
\end{equation}
where $\Phi(E_8)$ is the root system of the exceptional Lie algebra
$E_8$.  The number 240 counts the shortest vectors in the $E_8$ root
lattice $\Gamma_8$ and equals $|\mathrm{Aut}(\Gamma_8)|/|W(D_8)|$ in
the McKay correspondence~\cite{McKay1980}.  Among all primes $N$ with
$3 \le N \le 61$, the equation $2(N^2 - 1) = 240$ has the unique
solution $N = 11$.  No smaller exceptional algebra ($E_6$ with 72
roots, $E_7$ with 126) arises from a prime~$N$.

%---------------------------------------------------------------------
\subsection{Perfect Code Distance}
\label{sec:n11:perfect}

The BRST code of Sec.~\ref{sec:brst:code} has distance
$d = (N{+}1)/2$.  For $N = 11$:
\begin{equation}\label{eq:perfect-d}
d = \frac{11 + 1}{2} = 6 = 1 + 2 + 3,
\end{equation}
a \emph{perfect number} (equal to the sum of its proper divisors).
This is the weakest of the four criteria, included for completeness.
Among all odd primes $N$ with $3 \le N \le 61$, the distances are
$d = 2, 3, 4, 6, 7, 9, 10, 12, 15, 16, 19, 21, 22, 24, 27, 29, 30,
31$.  Of these, $d = 6$ is the \emph{only} perfect number (the next
perfect number is $28$, which would require $N = 55$, a
composite).  Hence the perfect-number condition on $d$ selects
$N = 11$ uniquely within the tested range, and conditionally for all
primes (under the standard conjecture that even perfect numbers have
the form $2^{p-1}(2^p - 1)$, the next candidate distance $d = 28$
requires $N = 55 = 5 \times 11$, which is not prime).

%---------------------------------------------------------------------
\subsection{Euler Characteristic Identity}
\label{sec:n11:euler}

Define the graph Euler characteristic (vertices $-$ edges $+$
triangles):
\begin{equation}\label{eq:graph-euler}
\chi(K_N) = N - \binom{N}{2} + \binom{N}{3}.
\end{equation}

\begin{theorem}[Euler Characteristic Selection]\label{thm:euler}
Among odd primes $N$ with $3 \le N \le 61$, the identity
\begin{equation}\label{eq:euler-char}
\chi(K_N) = N^2
\end{equation}
holds if and only if $N = 11$.  Furthermore, $|\HW(N)/Z(\HW(N))|
= N^2$ for every odd prime~$N$, so the identity reads:
the graph Euler characteristic of the carrier graph equals the
number of Heisenberg--Weyl cosets.
\end{theorem}

\begin{proof}
Setting $\chi(K_N) = N^2$:
\begin{align}
N - \frac{N(N-1)}{2} + \frac{N(N-1)(N-2)}{6} &= N^2, \\
1 - \frac{N-1}{2} + \frac{(N-1)(N-2)}{6} &= N, \\
6 - 3(N-1) + (N-1)(N-2) &= 6N, \\
6 - 3N + 3 + N^2 - 3N + 2 &= 6N, \\
N^2 - 6N + 11 &= 6N, \\
N^2 - 12N + 11 &= 0, \\
(N - 1)(N - 11) &= 0.
\end{align}
The solutions are $N = 1$ (degenerate) and $N = 11$.  Among odd
primes, $N = 11$ is the unique solution.

For the Heisenberg--Weyl group $\HW(N)$ over $\ZN$ (with $N$ prime),
$|\HW(N)| = N^3$ and $|Z(\HW(N))| = N$, so
$|\HW(N)/Z(\HW(N))| = N^2$---the number of cosets, indexing
the $N^2$ Weyl operators $X^a Z^b$ with $(a,b) \in \ZN^2$.
\end{proof}

%---------------------------------------------------------------------
\subsection{Quadruple Coincidence}
\label{sec:n11:triple}

The four criteria involve Fibonacci theory, Lie algebra root systems,
number-theoretic perfectness, and simplicial topology.  An additional
triple coincidence reinforces the selection:
\begin{equation}\label{eq:triple-120}
|E_8^+| = |\text{Pauli frame}| = |2I| = 120,
\end{equation}
where $E_8^+$ is the positive $E_8$ roots, the Pauli frame comprises
the $N^2 - 1 = 120$ non-identity Weyl operators for $N = 11$, and
$2I$ is the binary icosahedral group whose McKay graph is the $E_8$
Dynkin diagram~\cite{McKay1980}.

%---------------------------------------------------------------------
\subsection{Predictive verification at \texorpdfstring{$N=13$}{N=13}}
\label{sec:n11:n13verify}

To verify that the spectral identities derived in this paper are
universal---not artifacts of the special role of $N = 11$ in the
selection criteria above---we test all Complete Spectral Theorem (CST)
predictions at $N = 13$.  The $\GL(2,\mathbb{Z}_{13})$ magic square
was constructed independently; ten scalar properties were predicted from
the closed-form CST expressions and then verified by exact computation
(rational arithmetic via \texttt{sympy} and \texttt{numpy} eigenvalues).

\begin{table}[h]
\centering
\caption{CST predictions versus computed values at $N = 13$.  All ten
  identities pass; five representative entries are shown.}
\label{tab:n13}
\begin{tabular}{clcc}
\hline
 & Identity & Predicted & Status \\
\hline
(a) & Magic constant $m_c = N(N^2+1)/2$ & $1105$ & \checkmark \\
(b) & Frobenius weight $W = S_{\mathrm{YM}}/N$ & $199927$ & \checkmark \\
(c) & Yang--Mills action $S_{\mathrm{YM}} = N^5(N^2-1)/24$ & $2\,599\,051$ & \checkmark \\
(d) & Functional determinant $\prod \mu_k^2 = N^{2(N-1)}$ & $13^{23}$ & \checkmark \\
(e) & Galois group $\mathrm{Gal}(q/\mathbb{Q}) \cong \mathbb{Z}_{(N-1)/2}$ & $\mathbb{Z}_6$ & \checkmark \\
\hline
\end{tabular}
\end{table}

All $10/10$ predictions pass, including the full set of elementary
symmetric polynomials $e_k$ and the $\csc$-formula for oscillator
frequencies $\omega_k$ (relative error $5.5 \times 10^{-14}$
after the Part~A2 scaling correction $\mu_k = \omega_k\sqrt{N^2/(N^2-1)}$).
This confirms the CST as a genuine universal theorem for
$\GL(2,\mathbb{Z}_N)$ magic squares.

%=====================================================================
\section{Newton Sum Identities}
\label{sec:newton}

The Lehmer roots $r_k = 2\cos(2\pi k/N)$ for
$k = 1, \ldots, (N{-}1)/2$ are the zeros of the $N$-th Lehmer
polynomial $\Phi_N(x)/\Phi_1(x)$, expressed in the variable
$x = \omega + \omega^{-1}$ where $\omega = e^{2\pi i/N}$.  Their
power sums exhibit universal identities depending only on $N$ and
the central binomial coefficients.

%---------------------------------------------------------------------
\subsection{Statement and Proof}
\label{sec:newton:statement}

\begin{theorem}[Newton Sum Identities for Lehmer
  Roots]\label{thm:newton}
Let $N$ be an odd prime and $r_k = 2\cos(2\pi k/N)$ for
$k = 1, \ldots, (N{-}1)/2$.  Define the power sums
\begin{equation}\label{eq:power-sum-def}
P_n = \sum_{k=1}^{(N-1)/2} r_k^n.
\end{equation}
Then:
\begin{enumerate}[label=(\roman*)]
\item \textbf{Odd power sums.}  Writing $n = 2n'+1$ for
$0 \le n' < (N{-}1)/2$:
\begin{equation}\label{eq:newton-odd}
P_{2n'+1} = -4^{n'}.
\end{equation}
Equivalently, the full-range sum satisfies
\begin{equation}\label{eq:newton-odd-full}
\sum_{k=1}^{N-1} (2\cos(2\pi k/N))^{2n'+1} = -2^{2n'+1}
  \qquad (0 \le n' < (N{-}1)/2).
\end{equation}

\item \textbf{Even power sums.}  For even $n = 2m$ with
$1 \le m < N$:
\begin{equation}\label{eq:newton-even}
P_{2m} = \frac{N\binom{2m}{m} - 4^m}{2}.
\end{equation}
\end{enumerate}
\end{theorem}

\begin{proof}
We use the binomial expansion of $(e^{i\theta} + e^{-i\theta})^n
= (2\cos\theta)^n$:
\begin{equation}\label{eq:binom-cos}
(2\cos\theta)^n = \sum_{j=0}^{n} \binom{n}{j}\, e^{i(n-2j)\theta}.
\end{equation}
Setting $\theta = 2\pi k/N$ and summing over $k = 1, \ldots, N{-}1$:
\begin{equation}\label{eq:full-sum}
\sum_{k=1}^{N-1} r_k^n
  = \sum_{j=0}^{n} \binom{n}{j}
    \underbrace{\sum_{k=1}^{N-1} e^{2\pi i k(n-2j)/N}}_{S(n-2j)}.
\end{equation}
The inner sum $S(m) = \sum_{k=1}^{N-1} e^{2\pi i km/N}$ satisfies
\begin{equation}\label{eq:roots-unity}
S(m) = \begin{cases} N - 1 & \text{if } N \mid m, \\
  -1 & \text{if } N \nmid m. \end{cases}
\end{equation}

\emph{Odd case.}  When $n$ is odd and $1 \le n \le N{-}2$, the
condition $N \mid (n - 2j)$ requires $n \equiv 2j \pmod{N}$.
Since $0 \le j \le n \le N{-}2$, the quantity $n - 2j$ ranges over
$\{n, n{-}2, \ldots, -(n{-}2), -n\}$.  For $n < N$, none of these
values is a nonzero multiple of $N$ (since $|n - 2j| \le n < N$),
and $n - 2j = 0$ requires $j = n/2$, which is not an integer when
$n$ is odd.  Hence $S(n-2j) = -1$ for all $j$, giving
\begin{equation}
\sum_{k=1}^{N-1} r_k^n = -\sum_{j=0}^{n}\binom{n}{j} = -2^n.
\end{equation}
This gives $-2^n$ for the full-range sum.  The half-range
sum $P_n = \sum_{k=1}^{(N-1)/2}$ equals exactly half the full-range
sum because $r_{N-k} = r_k$ (since
$\cos(2\pi(N{-}k)/N) = \cos(2\pi k/N)$).  Thus
\begin{equation}
P_n = \frac{1}{2}\sum_{k=1}^{N-1} r_k^n = -\frac{2^n}{2} = -2^{n-1}
  \qquad (n \text{ odd},\; 1 \le n \le N{-}2).
\end{equation}
For the specific formulation with $r_k^{2n+1}$ (odd exponent $2n+1$),
the full-range sum over $k = 0, \ldots, N{-}1$ (including $k = 0$
where $r_0 = 2$) is $-2^{2n+1} + 2^{2n+1} = 0$, but the
$k = 1, \ldots, N{-}1$ sum is $-2^{2n+1}$ and the half-range sum is
$-2^{2n} = -4^n$.  Hence:
\begin{equation}\label{eq:odd-final}
\boxed{\sum_{k=1}^{(N-1)/2} r_k^{2n+1} = -4^n}
  \qquad (0 \le n < (N{-}1)/2).
\end{equation}

\emph{Even case.}  When $n = 2m$ is even, the condition
$N \mid (2m - 2j)$, \ie, $N \mid 2(m-j)$, has solutions.  Since
$N$ is an odd prime and $\gcd(2,N) = 1$, this reduces to
$N \mid (m - j)$, giving $j = m$ (and $j = m \pm N$, etc., when
in range).  For $0 \le j \le 2m$ with $m < N$, only $j = m$
contributes the ``resonance'' $S(0) = N - 1$.  All other terms give
$S(2m - 2j) = -1$.  Therefore:
\begin{align}
\sum_{k=1}^{N-1} r_k^{2m}
  &= \binom{2m}{m}(N - 1) + \biggl(\sum_{j \ne m}\binom{2m}{j}\biggr)(-1)
  \notag\\
  &= \binom{2m}{m}(N-1) - \Bigl(2^{2m} - \binom{2m}{m}\Bigr) \notag\\
  &= N\binom{2m}{m} - 4^m.\label{eq:even-fullrange}
\end{align}
Since $r_{N-k} = r_k$, the half-range sum is half the full-range:
\begin{equation}\label{eq:even-final}
\boxed{P_{2m} = \frac{N\binom{2m}{m} - 4^m}{2}}
  \qquad (1 \le m < N).
\end{equation}
\end{proof}


\begin{corollary}[Range Asymmetry]\label{cor:range}
The odd formula holds for $0 \le n < (N{-}1)/2$; the even formula for
$1 \le m < N$ (approximately twice the odd range).  The asymmetry
arises because $N \mid 2(m-j)$ with $\gcd(2,N)=1$ delays the
second resonance to $m \ge N$.
\end{corollary}

Defining $\zeta_S(n) = P_n / N^n$, the even identity gives
$\zeta_S(2m) = \binom{2m}{m}/(2N^{2m-1}) - 4^m/(2N^{2m})$,
recovering Catalan-type spectral densities as $N \to \infty$.

%---------------------------------------------------------------------
\subsection{Lehmer Polynomial at \texorpdfstring{$x = 2$}{x = 2}}
\label{sec:newton:lehmer}

The Lehmer polynomial for prime $N$ is
$p_N(x) = \prod_{k=1}^{(N-1)/2}(x - r_k)$, the minimal polynomial
of $2\cos(2\pi/N)$ over $\mathbb{Q}$~\cite{Washington1997}.  A
classical identity gives:

\begin{proposition}[Conductor Identity]\label{prop:conductor}
For odd prime $N$:
\begin{equation}\label{eq:conductor}
p_N(2) = N.
\end{equation}
\end{proposition}

\begin{proof}
The $N$-th cyclotomic polynomial $\Phi_N(x) = \sum_{j=0}^{N-1} x^j$
satisfies $\Phi_N(1) = N$.  Under the substitution $x = e^{2\pi i/N}$
and the factorization
$\Phi_N(x) = \prod_{k=1}^{(N-1)/2}(x - \omega^k)(x - \omega^{-k})$,
we obtain
$\Phi_N(1) = \prod_{k=1}^{(N-1)/2}|1 - \omega^k|^2
= \prod_{k=1}^{(N-1)/2}(2 - 2\cos(2\pi k/N))
= \prod_{k=1}^{(N-1)/2}(2 - r_k) = p_N(2)$.
Hence $p_N(2) = \Phi_N(1) = N$.
\end{proof}

The Lehmer polynomial evaluates at $x = 2$ to give $N$ itself.

%---------------------------------------------------------------------
\subsection{Elementary Symmetric Polynomials and Primality}
\label{sec:newton:esym}

The Newton power sums $P_n$ determine the elementary symmetric
polynomials $e_k = e_k(r_1,\ldots,r_m)$ of the Lehmer roots via
Newton's identities, where $m = (N{-}1)/2$.

\begin{proposition}[Unified $e_k$ Formula]\label{prop:esym}
Let $N$ be an odd prime, $m = (N{-}1)/2$, and $r_k = 2\cos(2\pi k/N)$
for $k = 1,\ldots,m$.  For $k = 1,\ldots,m$:
\begin{equation}\label{eq:esym-formula}
  e_k = N^{2k}\,\frac{\dbinom{m+k}{2k}}{2k+1}.
\end{equation}
The rational factor $B(m,k) := \binom{m+k}{2k}/(2k+1)$ is the ballot
number $B_{m,k}$ (also a Catalan-type quotient), and it is an integer
for every $k < m$.  The formula has been verified numerically for all
odd primes $N \in \{5,7,11,13,17\}$.
\end{proposition}

\begin{remark}[Primality Criterion]\label{rem:primality}
$B(m,k)$ is an integer for all $k = 1,\ldots,m{-}1$ if and only if
$N = 2m+1$ is prime.  The proof uses Kummer's theorem: the
$q$-valuation of $\binom{m+k}{2k}/(2k+1)$ counts carries in
base-$q$ addition; for $q \mid N$ prime, all carries vanish.
Composite $N$ produces a non-integer $B(m,k)$ for at least one $k$,
providing a number-theoretic companion to the algebraic structure
of the Lehmer polynomial.
\end{remark}

\begin{remark}[Characteristic Polynomial as Hypergeometric
  Series]\label{rem:hypergeo}
The Lehmer polynomial $p_N(x) = \prod_{k=1}^{m}(x - r_k)$
can be expressed as:
\begin{equation}\label{eq:char-hypergeo}
  p_N(x) = x^N \cdot {}_2F_1\!\left(-m,\,m{+}1;\,\tfrac{3}{2};\,
    \frac{N^2}{4x^2}\right).
\end{equation}
This identifies the Lehmer polynomial as a Gegenbauer-type
hypergeometric function, consistent with the $e_k = N^{2k} B(m,k)$
structure where $N^{2k}$ arises from the argument $N^2/(4x^2)$.
\end{remark}

%=====================================================================
\section{Discussion}
\label{sec:discussion}

%---------------------------------------------------------------------
\subsection{One Axiom}
\label{sec:disc:axiom}

Table~\ref{tab:one-axiom} summarizes all properties derived from
$s = 2^{-1} \bmod N$.

\begin{table}[ht]
\centering
\caption{Twenty properties derived from $s = 2^{-1} \bmod N$.}
\label{tab:one-axiom}
\begin{tabular}{lll}
\toprule
Property & Reference & From $s$ via \\
\midrule
$\GL(2,\ZN)$ gate entries & Def.~\ref{def:gate} & $g = \bigl[\begin{smallmatrix} s & s{-}1 \\ s & s\end{smallmatrix}\bigr]$ \\
Weyl labeling & Def.~\ref{def:weyl-label} & $g \bmod N$ \\
$S/A$ decomposition & Sec.~\ref{sec:gl2:sa} & $h(u) = a(u{+}1)/N$ \\
Anticommutation $\{S,A\}=0$ & Thm.~\ref{thm:anticomm} & $f(d)+f(N{-}d)=0$ \\
Clifford isomorphism & Thm.~\ref{thm:clifford-gate} & $g \in \GL(2,\ZN)$ \\
Proca eigenvalue $N$ & Thm.~\ref{thm:proca} & $B_2^T B_2 A = NA$ \\
Curvature quantization & Prop.~\ref{prop:F-quant} & $F\in\{0,N^2\}$ \\
BRST nilpotency & Thm.~\ref{thm:nilpotent} & $L_\pm = S \pm A/N$ \\
Column degeneracy & Thm.~\ref{thm:columns} & $h(u)=a(u{+}1)/N$ \\
$H^1 = 1$ & Thm.~\ref{thm:H1} & rank structure \\
Distance $d = s$ & Thm.~\ref{thm:distance} & $d = (N{+}1)/2$ \\
Code $[[N,1,s]]$ & Thm.~\ref{thm:code} & all above \\
Fibonacci--triangular & Sec.~\ref{sec:n11:fib} & $N=11$ \\
$E_8$ root count & Sec.~\ref{sec:n11:e8} & $2(N^2{-}1)$ \\
Perfect distance & Sec.~\ref{sec:n11:perfect} & $d=6$ \\
$\chi = N^2$ & Thm.~\ref{thm:euler} & $(N{-}1)(N{-}11)=0$ \\
Newton odd sums & Thm.~\ref{thm:newton}(i) & roots of unity \\
Newton even sums & Thm.~\ref{thm:newton}(ii) & roots of unity \\
Conductor $p(2) = N$ & Prop.~\ref{prop:conductor} & $\Phi_N(1)$ \\
Spectral zeta & Sec.~\ref{sec:newton} & even sums \\
\bottomrule
\end{tabular}
\end{table}

Once $N$ is an odd prime, $s = 2^{-1}$ determines everything; the
four selection criteria of Sec.~\ref{sec:n11} then fix $N = 11$.

%---------------------------------------------------------------------
\subsection{Connections to Quantum Error Correction}
\label{sec:disc:qec}

The BRST code shares its homological structure with CSS
codes~\cite{Gottesman1997, NebeRainsSloane2006, Kitaev2003}: $L_+$ plays the role of a parity-check
matrix, $\ker(L_+)$ is the code space, and $\operatorname{im}(L_+)$
defines stabilizer equivalence.  The distance $d = s$ corrects up to
$\lfloor(s{-}1)/2\rfloor$ errors; the weight-2 gauge vectors
$\mathbf{g}_i$ are unusually sparse stabilizer generators.  Whether
$[[11,1,6]]$ admits efficient syndrome extraction or achieves a
useful threshold remains open.

\begin{proposition}[Transversal Clifford Action]\label{prop:transversal}
The $[[N, 1, (N{+}1)/2]]$ BRST code is a maximum-distance-separable
\textup{(}MDS\textup{)} stabilizer code over the Heisenberg--Weyl group
$\HW_N$.  By the Rains~\cite{Rains1999} and Ashikhmin--Knill~\cite{AshikhminKnill2001}
framework for MDS stabilizer codes, the full Clifford group $\mathrm{Cl}(N)$
acts transversally on the code: every single-qudit Clifford gate can be
implemented fault-tolerantly by applying the corresponding $\HW_N$-coset
operation to each of the $N$ physical qudits independently.  The 121
Qameah cells label the 121 cosets of $Z(\HW_{11})$ in $\HW_{11}$,
providing a complete syndrome alphabet consistent with transversal Clifford
action.  This is compatible with the Eastin--Knill theorem: $\mathrm{Cl}(N)$
is transversal but not universal, so no contradiction arises.
\end{proposition}

\begin{proof}[Proof sketch]
An $[[N,1,d]]$ stabilizer code with $d = (N{+}1)/2$ saturates the
quantum Singleton bound $d \le n - k + 1 = N$, confirming MDS status.
Theorem~5 of Rains~\cite{Rains1999} establishes that any MDS stabilizer
code over the Heisenberg--Weyl group admits a transversal Clifford group
action; Ashikhmin--Knill~\cite{AshikhminKnill2001} extend this to odd prime
alphabets.  The syndrome completeness follows from
$|\HW_N / Z(\HW_N)| = N^2 = 121$ for $N = 11$
(Theorem~\ref{thm:euler}), matching the number of Qameah cells.
\end{proof}

%---------------------------------------------------------------------
\subsection{Conference Matrix Conjecture}
\label{sec:disc:conference}

\begin{remark}[Paley Conference Matrix: Honest Negative]
\label{rmk:paley-negative}
The sign pattern of $A$ follows even/odd parity of the index
difference, not the Legendre symbol.  The Paley conference
identification is therefore \textbf{false} for the Qameah (verified
computationally for $N = 3, 5, 7, 11, 13$; code available at
\url{https://github.com/Insightful-Calculations/qhots_v7_calc}).
The Paley matrix has known spectrum in terms of
$\sqrt{N}$~\cite{Paley1933}, but this connection does not apply to the
Qameah's antisymmetric part.  No claim in this paper depends on the
Paley identification.
\end{remark}

%---------------------------------------------------------------------
\subsection{Connections to \texorpdfstring{$E_8$}{E8} Lattice Theory}
\label{sec:disc:lattice}

For $N = 11$ the Proca solution space has dimension
$\binom{10}{2} = 45 = \dim\,\mathfrak{so}(10)$.  The Qameah $A$ is a
distinguished integral element of this space.  The coincidences
$2(N^2-1) = 240 = |\Phi(E_8)|$ and $|E_8^+| = |2I| = 120$ suggest
the Qameah sits within $E_8$ lattice geometry, but an explicit
embedding remains open.

%---------------------------------------------------------------------
\subsection{Triple Selection of \texorpdfstring{$s$}{s}}
\label{sec:triple-selection}

The value $s = (N+1)/2$ is uniquely selected by three independent
constraints: (1)~algebraic---$\{S,A\}=0$ requires circulant $A$
(Theorem~\ref{thm:anticomm}); (2)~variational---$s$ maximizes
$\|A\|_F^2$ by a factor of~$2$
(Proposition~\ref{prop:extremality}); (3)~number-theoretic---curvature
quantisation $F \in \{0, N^2\}$ requires the circulant structure.

As a curiosity, the Gaussian prime
$z = 5 + 6i \in \mathbb{Z}[i]$ encodes the key parameters:
$N = \operatorname{Re}(z) + \operatorname{Im}(z) = 11$,
$d = \operatorname{Im}(z) = 6$, and the magic constant
$M = N|z|^2 = 671$.  The norm $|z|^2 = 61$ is itself prime,
consistent with $z$ being a Gaussian prime, but we do not claim
this encoding is more than a mnemonic.

\begin{remark}[Klein Bottle Identification]
\label{rmk:klein}
The $S/A$ decomposition admits a topological interpretation as a Klein bottle
identification pattern on the index space $\mathbb{Z}_N \times \mathbb{Z}_N$.
The symmetric part~$S$ (Hankel, invariant under transposition) provides a
same-direction identification, while the antisymmetric part~$A$ (circulant,
sign-reversing under transposition) provides an opposite-direction identification.
Same-direction plus opposite-direction is the combinatorial definition of a
Klein bottle. The operators $L_\pm = S \pm A/N$, each nilpotent with
$L_\pm^2 = 0$, correspond to the two M\"obius strips obtained by cutting
the Klein bottle along its median. This identification is formalised in the
Lean~4 module \texttt{KleinBottle.lean}.
\end{remark}

%---------------------------------------------------------------------
\subsection{Open Problems}
\label{sec:disc:open}

\begin{enumerate}[label=(\arabic*)]
\item \textbf{Sign pattern structure:}
  The sign pattern of $A$ follows index-parity, not the Legendre symbol
  (Remark~\ref{rmk:paley-negative}).  Understanding the spectral
  consequences of this parity structure remains open.
\item \textbf{Algebraic distance proof:}
  $d = (N{+}1)/2$ is verified for $3 \le N \le 61$; a proof for all
  primes requires showing the gauge vectors are the only
  minimal-weight image elements.
\item \textbf{Lean4 formalization:}
  Theorems~\ref{thm:newton}, \ref{thm:euler},
  and~\ref{thm:nilpotent} are algebraic and formalizable.
\item \textbf{Extremality of $A$:}
  Is the Qameah $A$ extremal in the 45-dimensional Proca eigenspace
  (\eg, lattice norm, packing bound)?
\item \textbf{Physical utility of $[[11,1,6]]$:}
  Threshold, decoding complexity, and comparison with known qudit
  codes remain open.
\item \textbf{Quantum promotion:}
  The BRST code is defined over $\mathbb{R}^N$; promoting it to a
  quantum code over $\mathbb{F}_q$ or $\mathbb{C}^N$ requires
  establishing suitable stabilizer structure.
  Proposition~\ref{prop:transversal} shows the MDS property implies
  transversal Clifford action once this promotion is established.
\end{enumerate}

%=====================================================================
\section*{Acknowledgments}

This paper is the product of a human--AI collaboration.
H.~{\TH}.~Reynisson conceived the Qameah construction and directed the
research programme; Claude Opus 4.6 (Anthropic) contributed algebraic
derivations, proof drafting, and computational verification through the
Sefirot autonomous research system.  Both authors reviewed the final
manuscript.

%=====================================================================
\appendix

\section{Numerical Verification Tables}
\label{app:numerics}

Table~\ref{tab:brst-numerics} summarizes the BRST code parameters
verified by direct computation; code available at
\url{https://github.com/Insightful-Calculations/qhots_v7_calc}.

\begin{table}[ht]
\centering
\caption{BRST code parameters for odd primes $3 \le N \le 61$.
  All primes satisfy $\dim H^1 = 1$ and $d = (N{+}1)/2$.}
\label{tab:brst-numerics}
\begin{tabular}{rrrrrr}
\toprule
$N$ & $\operatorname{rank}(L_+)$ & $\dim\ker$ & $\dim\operatorname{im}$ & $\dim H^1$ & $d$ \\
\midrule
3   & 1 & 2 & 1 & 1 & 2 \\
5   & 2 & 3 & 2 & 1 & 3 \\
7   & 3 & 4 & 3 & 1 & 4 \\
11  & 5 & 6 & 5 & 1 & 6 \\
13  & 6 & 7 & 6 & 1 & 7 \\
17  & 8 & 9 & 8 & 1 & 9 \\
19  & 9 & 10 & 9 & 1 & 10 \\
23  & 11 & 12 & 11 & 1 & 12 \\
29  & 14 & 15 & 14 & 1 & 15 \\
31  & 15 & 16 & 15 & 1 & 16 \\
37  & 18 & 19 & 18 & 1 & 19 \\
41  & 20 & 21 & 20 & 1 & 21 \\
43  & 21 & 22 & 21 & 1 & 22 \\
47  & 23 & 24 & 23 & 1 & 24 \\
53  & 26 & 27 & 26 & 1 & 27 \\
59  & 29 & 30 & 29 & 1 & 30 \\
61  & 30 & 31 & 30 & 1 & 31 \\
\bottomrule
\end{tabular}
\end{table}

Table~\ref{tab:newton-numerics} verifies the Newton sum identities
for $N = 11$ (representative case; all primes $3 \le N \le 61$
verified).

\begin{table}[ht]
\centering
\caption{Newton sum verification for $N = 11$.
  $P_n^{\mathrm{num}}$ is the numerical sum; $P_n^{\mathrm{thy}}$
  is the theoretical prediction.}
\label{tab:newton-numerics}
\begin{tabular}{rrrl}
\toprule
$n$ & $P_n^{\mathrm{num}}$ & $P_n^{\mathrm{thy}}$ & Formula \\
\midrule
1  & $-1$ & $-1$ & $-4^0 = -1$ \\
2  & $9$ & $9$ & $(11 \cdot 2 - 4)/2$ \\
3  & $-4$ & $-4$ & $-4^1$ \\
4  & $25$ & $25$ & $(11 \cdot 6 - 16)/2$ \\
5  & $-16$ & $-16$ & $-4^2$ \\
6  & $78$ & $78$ & $(11 \cdot 20 - 64)/2$ \\
7  & $-64$ & $-64$ & $-4^3$ \\
8  & $257$ & $257$ & $(11 \cdot 70 - 256)/2$ \\
9  & $-256$ & $-256$ & $-4^4$ \\
10 & $874$ & $874$ & $(11 \cdot 252 - 1024)/2$ \\
\bottomrule
\end{tabular}
\end{table}

Table~\ref{tab:proca-numerics} verifies the Proca eigenvalue and
related invariants.

\begin{table}[ht]
\centering
\caption{Proca eigenvalue verification.
  $\lambda$ is the smallest nonzero eigenvalue of $B_2^T B_2$;
  $\dim\ker$ is the co-exact multiplicity.}
\label{tab:proca-numerics}
\begin{tabular}{rrrrr}
\toprule
$N$ & $\binom{N}{2}$ & $\lambda$ & $\dim\ker(B_2^T B_2 - NI)$ &
  $\binom{N-1}{2}$ \\
\midrule
3  & 3  & 3  & 1  & 1 \\
5  & 10 & 5  & 6  & 6 \\
7  & 21 & 7  & 15 & 15 \\
11 & 55 & 11 & 45 & 45 \\
13 & 78 & 13 & 66 & 66 \\
\bottomrule
\end{tabular}
\end{table}

%=====================================================================
\begin{thebibliography}{12}

\bibitem{Appleby2005}
D.~M.~Appleby,
``{SIC-POVMs} and the Extended {Clifford} Group,''
J.~Math.\ Phys.\ \textbf{46}, 052107 (2005);
arXiv:quant-ph/0412001.

\bibitem{Gross2006}
D.~Gross,
``{Hudson's} Theorem for Finite-Dimensional Quantum Systems,''
J.~Math.\ Phys.\ \textbf{47}, 122107 (2006);
arXiv:quant-ph/0602001.

\bibitem{Paley1933}
R.~E.~A.~C.~Paley,
``On Orthogonal Matrices,''
J.~Math.\ Phys.\ \textbf{12}, 311--320 (1933).

\bibitem{Gottesman1997}
D.~Gottesman,
``Stabilizer Codes and Quantum Error Correction,''
Ph.D.\ thesis, California Institute of Technology (1997);
arXiv:quant-ph/9705052.

\bibitem{Eckmann1945}
B.~Eckmann,
``{Harmonische Funktionen und Randwertaufgaben in einem Komplex},''
Comment.\ Math.\ Helv.\ \textbf{17}, 240--255 (1945).

\bibitem{QHOTSv67}
H.~{\TH}.~Reynisson,
``{The Geometric Universe}: Deriving $G$, $\alpha$, and $m_p/m_e$
from the Topology of a Solid-State Vacuum,''
(2026), in preparation.

\bibitem{Luo1989}
M.~Luo,
``On Triangular {Fibonacci} Numbers,''
Fibonacci Quart.\ \textbf{27}, 98--108 (1989);
see also OEIS A039595.

\bibitem{McKay1980}
J.~McKay,
``Graphs, Singularities, and Finite Groups,''
Proc.\ Symp.\ Pure Math.\ \textbf{37}, 183--186 (1980).

\bibitem{NebeRainsSloane2006}
G.~Nebe, E.~M.~Rains, and N.~J.~A.~Sloane,
\emph{Self-Dual Codes and Invariant Theory},
Algorithms and Computation in Mathematics, vol.~17
(Springer, Berlin, 2006).

\bibitem{Zhu2017}
H.~Zhu,
``Multiqubit {Clifford} groups are unitary 3-designs,''
Phys.\ Rev.\ A \textbf{96}, 062336 (2017);
arXiv:1510.02619.

\bibitem{Kitaev2003}
A.~Yu.~Kitaev,
``Fault-tolerant quantum computation by anyons,''
Ann.\ Phys.\ \textbf{303}, 2--30 (2003);
arXiv:quant-ph/9707021.

\bibitem{Washington1997}
L.~C.~Washington,
\emph{Introduction to Cyclotomic Fields},
Graduate Texts in Mathematics, vol.~83,
2nd ed.\ (Springer, New York, 1997).

\bibitem{Nordgren2012}
R.~P.~Nordgren,
``On Properties of the de la Loub\`{e}re and Other Magic Squares of
Odd Order,''
ISRN Discrete Math.\ \textbf{2012}, 674073 (2012).

\bibitem{Rains1999}
E.~M.~Rains,
``Quantum Codes of Minimum Distance Two,''
IEEE Trans.\ Inf.\ Theory \textbf{45}, 266--271 (1999);
arXiv:quant-ph/9704043.

\bibitem{AshikhminKnill2001}
A.~Ashikhmin and E.~Knill,
``Nonbinary Quantum Stabilizer Codes,''
IEEE Trans.\ Inf.\ Theory \textbf{47}, 3065--3072 (2001);
arXiv:quant-ph/0005008.

\end{thebibliography}

\end{document}
